% concatenated from tang_arcs_fol_DV2.tex
\documentclass[11pt]{amsart}

% Encoding and language
\usepackage[british]{babel}
\usepackage[T1]{fontenc}
\usepackage[utf8]{inputenc}

% Mathematics
\usepackage{amsmath,amssymb,amsfonts,amsthm,mathtools}
\usepackage{mathrsfs}
\usepackage{latexsym}
\usepackage{textcomp}

% Layout and tables
\usepackage{array}
\usepackage{booktabs}
\usepackage{enumitem}
\usepackage{setspace}

% Diagrams and graphics
\usepackage{tikz}
\usepackage{tikz-cd}
\usetikzlibrary{trees}
\usepackage{amscd}

% Fonts
\usepackage{times}
\usepackage{microtype}

% Hyperlinks
\usepackage{hyperref}
\hypersetup{hidelinks}

% Page layout matching the old style
\setlength{\textheight}{210mm}
\setlength{\topmargin}{0.46cm}
\setlength{\textwidth}{152mm}
\setlength{\evensidemargin}{0.60cm}
\setlength{\oddsidemargin}{0.60cm}

\setcounter{MaxMatrixCols}{30}
\numberwithin{equation}{section}
\setcounter{tocdepth}{1}
\renewcommand{\sectionmark}[1]{\markright{\thesection \ #1}{}}

% Theorem environments
\theoremstyle{plain}
\newtheorem{Theorem}{Theorem}[section]
\newtheorem{theoremn}{Theorem}
\newtheorem{Lemma}[Theorem]{Lemma}
\newtheorem{Proposition}[Theorem]{Proposition}
\newtheorem{Corollary}[Theorem]{Corollary}
\newtheorem{Conjecture}[Theorem]{Conjecture}
\newtheorem*{question}{Question}
\newtheorem*{remark}{Remark}
\newtheorem*{conjecturen}{Conjecture}

\theoremstyle{definition}
\newtheorem{Definition}[Theorem]{Definition}
\newtheorem{Construction}[Theorem]{Construction}
\newtheorem{Notation}[Theorem]{Notation}
\newtheorem{Example}[Theorem]{Example}
\newtheorem{Remark}[Theorem]{Remark}
\newtheorem{Question}[Theorem]{Question}
\newtheorem{Assumption}[Theorem]{Assumption}
\newtheorem{say}[Theorem]{}

% General macros
\newcommand{\arXiv}[1]{\href{https://arxiv.org/abs/#1}{arXiv:#1}}
\newcommand{\ndo}{\node[draw,circle,inner sep=1]}

\newcommand{\Spec}{\operatorname{Spec}}
\newcommand{\Hom}{\operatorname{Hom}}
\newcommand{\Spf}{\operatorname{Spf}}
\newcommand{\Proj}{\operatorname{Proj}}
\newcommand{\Sing}{\operatorname{Sing}}
\newcommand{\Supp}{\operatorname{Supp}}
\newcommand{\Exc}{\operatorname{Exc}}
\newcommand{\codim}{\operatorname{codim}}
\newcommand{\ord}{\operatorname{ord}}
\newcommand{\val}{\operatorname{val}}
\newcommand{\mld}{\operatorname{mld}}
\newcommand{\tmldtor}{\operatorname{tmld}_{\mathrm{tor}}}
\newcommand{\lct}{\operatorname{lct}}
\newcommand{\tlct}{\operatorname{tlct}}
\newcommand{\lcodim}{\operatorname{lcodim}}
\newcommand{\Diff}{\operatorname{Diff}}
\newcommand{\Cont}{\operatorname{Cont}}
\newcommand{\Nklt}{\operatorname{Nklt}}
\newcommand{\Nlc}{\operatorname{Nlc}}
\newcommand{\TNklt}{\operatorname{Nklt}_{\tan}}
\newcommand{\Jac}{\operatorname{Jac}}
\newcommand{\Nash}{\operatorname{Nash}}
\newcommand{\coeff}{\operatorname{coeff}}
\newcommand{\Cone}{\operatorname{Cone}}
\newcommand{\Res}{\operatorname{Res}}
\newcommand{\Aut}{\operatorname{Aut}}
\newcommand{\Ext}{\operatorname{Ext}}
\newcommand{\Hilb}{\operatorname{Hilb}}
\newcommand{\Ker}{\operatorname{Ker}}
\newcommand{\Isom}{\operatorname{Isom}}
\newcommand{\im}{\operatorname{Im}}
\newcommand{\Sym}{\operatorname{Sym}}
\newcommand{\Pic}{\operatorname{Pic}}
\newcommand{\Bs}{\operatorname{Bs}}
\newcommand{\Bl}{\operatorname{Bl}}
\newcommand{\Gr}{\operatorname{Gr}}
\newcommand{\End}{\operatorname{End}}

% Project-specific operators and symbols
\newcommand{\MJ}{\mathrm{MJ}}
\newcommand{\sn}{\mathrm{sn}}
\newcommand{\tanMJ}{\mathrm{tan\text{-}MJ}}
\newcommand{\Sigmatan}{\Sigma_{\mathrm{tan}}}
\newcommand{\bbound}{\mathfrak b}
\newcommand{\Adm}{\mathsf{Adm}}
\newcommand{\SN}{\mathsf{SN}}
\newcommand{\red}{\mathrm{red}}
\newcommand{\Strata}{\mathsf{Str}}
\newcommand{\Cyl}{\mathsf{Cyl}}
\newcommand{\inv}{\mathrm{inv}}
\newcommand{\tmld}{\operatorname{tmld}}
\newcommand{\Jetm}{J_m}
\newcommand{\dd}{\mathrm d}
\newcommand{\eps}{\varepsilon}
\newcommand{\lra}{\longrightarrow}
\newcommand{\Rlambda}{\mathcal R_\lambda}
\newcommand{\AF}{A_{\mathcal F}}
\newcommand{\Dinv}{D_{\mathrm{inv}}}

% Blackboard letters
\newcommand{\bA}{\mathbb A}
\newcommand{\bC}{\mathbb C}
\newcommand{\bN}{\mathbb N}
\newcommand{\bQ}{\mathbb Q}
\newcommand{\bR}{\mathbb R}
\newcommand{\bZ}{\mathbb Z}
\newcommand{\PP}{\mathbb P}
\newcommand{\CC}{\mathbb C}
\newcommand{\ZZ}{\mathbb Z}

% Caligraphic letters
\newcommand{\cF}{\mathcal F}
\newcommand{\cG}{\mathcal G}
\newcommand{\cO}{\mathcal O}
\newcommand{\cB}{\mathcal B}
\newcommand{\cC}{\mathcal C}
\newcommand{\cJ}{\mathcal J}
\newcommand{\cN}{\mathcal N}
\newcommand{\cS}{\mathcal S}
\newcommand{\cT}{\mathcal T}

% Script letters
\newcommand{\sA}{\mathscr A}
\newcommand{\sB}{\mathscr B}
\newcommand{\sC}{\mathscr C}
\newcommand{\sD}{\mathscr D}
\newcommand{\sE}{\mathscr E}
\newcommand{\sF}{\mathscr F}
\newcommand{\sG}{\mathscr G}
\newcommand{\sH}{\mathscr H}
\newcommand{\sI}{\mathscr I}
\newcommand{\sJ}{\mathscr J}
\newcommand{\sK}{\mathscr K}
\newcommand{\sL}{\mathscr L}
\newcommand{\sM}{\mathscr M}
\newcommand{\sN}{\mathscr N}
\newcommand{\sO}{\mathscr O}
\newcommand{\sP}{\mathscr P}
\newcommand{\sQ}{\mathscr Q}
\newcommand{\sR}{\mathscr R}
\newcommand{\sS}{\mathscr S}
\newcommand{\sT}{\mathscr T}
\newcommand{\sU}{\mathscr U}
\newcommand{\sV}{\mathscr V}
\newcommand{\sW}{\mathscr W}
\newcommand{\sX}{\mathscr X}
\newcommand{\sY}{\mathscr Y}
\newcommand{\sZ}{\mathscr Z}

% Other symbols
\newcommand{\GG}{\mathfrak G}
\newcommand{\jaco}{\mathfrak j}
\newcommand{\OO}{\mathcal O}
\newcommand{\Poiss}{\operatorname{Poisson}}
\newcommand{\Fol}{\operatorname{Fol}}
\newcommand{\ad}{\operatorname{ad}}
\newcommand{\rk}{\operatorname{rk}}
\newcommand{\slthree}{\mathfrak{sl}_3}
\newcommand{\sltwo}{\mathfrak{sl}_2}
\newcommand{\Spanop}{\operatorname{Span}}
\newcommand\Span[1]{\langle#1\rangle}
\newcommand{\Rat}{\textup{RatCurves}}

\DeclareMathOperator{\rank}{rank}

\def\bibaut#1{{\sc #1}}

\title[Foliated and Mather--Jacobian discrepancies]{Foliated and Mather--Jacobian discrepancies via tangential arcs}

\author[M. Corr\^ea]{Maur\'icio Corr\^ea}

\address[Maur\'icio Corr\^ea]{Dipartimento di Matematica, Universit\`a degli Studi di Bari Aldo Moro, Via E. Orabona 4, I-70125 Bari, Italy}

\email{mauricio.barros@uniba.it}

\begin{document}


\begin{abstract}
This article develops a tangential arc-space approach to foliated discrepancies
for logarithmic simple co-rank one foliations on threefolds, relative to a
fixed invariant normal crossing separatrix divisor.  In the non-resonant
logarithmic case, reduced tangential arcs centred on the prescribed
tangential locus are confined to this divisor.  The tangential sector is
therefore presented, at the reduced arc level, by the normalised
separatrix--conductor system.  Foliated adjunction transfers the discrepancy
calculus to ordinary log pairs obtained by adjunction on the normalised
branches and conductors.  The arc-space theorem of
Ein--Musta\c{t}\u{a}--Yasuda, applied on these strata, then gives a
tangential codimension formula identifying logarithmic codimensions of
toroidal tangential divisorial cylinders with the corresponding tangential
discrepancies.  For toroidal invariant divisors read on branches, this
tangential discrepancy agrees with the usual foliated discrepancy.  The
resulting theory gives toroidal tangential inversion of adjunction, a
branch--conductor description of the tangential non-lc and non-klt loci, a
cylinder criterion for tangential log canonicity, lower semicontinuity of the
toroidal tangential minimal log discrepancy, and a relative Mather--Jacobian
refinement for the canonical image separatrix system.
\end{abstract}


\maketitle
%\tableofcontents



\section{Introduction}

Arc spaces furnish a natural bridge between birational geometry and families
of formal curves.  The theorem of Ein--Musta\c{t}\u{a}--Yasuda expresses log
discrepancies as codimensions of divisorial cylinders in arc spaces
\cite{EMY}.  Beyond the locally complete intersection case, this relation is
replaced by Mather--Jacobian theory and its arc-space interpretation, due to
Ishii, Ein--Ishii, and de Fernex--Docampo \cite{IshiiMather,EI,dFD}.

We construct a foliated tangential analogue of this picture for co-rank one
foliations on threefolds.  The birational background is the foliated minimal
model programme, especially the role of simple non-dicritical singularities,
separatrices, and adjunction on invariant divisors
\cite{CSMMP,Spicer,SS,CSAdj}.  We work in the logarithmic simple adapted
non-resonant setting, relative to an invariant normal crossing separatrix
divisor \(\Dinv\).  The positive non-resonance condition has a precise role:
it prevents tangential arcs centred on \(\Dinv\) from leaving this divisor.
Canonical or log canonical singularities alone do not provide this local
arc-confinement statement.

The main local result is the reduced tangential arc-confinement theorem.  In
adapted logarithmic coordinates, if a tangential arc centred on \(\Dinv\) is
not contained in \(\Dinv\), then the lowest-order term of the pulled-back
logarithmic form gives a positive resonance among the logarithmic residues.
Thus, in the non-resonant sector, reduced tangential arcs are identified with
arcs on \(\Dinv\).  Passing to the seminormal branch--conductor system
then presents the reduced tangential sector by the ordinary arc spaces of
normalised branches and conductors.  The construction is deliberately made at
the level of reduced infinite arcs, not at the level of all finite Carter
tangent jet schemes \cite{Carter,CarterThesis}.

Foliated adjunction transfers the discrepancy calculus to ordinary log pairs
on the normalised branch--conductor strata.  These log pairs are obtained by
normalised adjunction.  This gives a tangential
Ein--Musta\c{t}\u{a}--Yasuda type formula for adapted toroidal tangential
divisorial cylinders.  For adapted toroidal invariant divisors whose data are
read on normalised branches, the resulting tangential discrepancy agrees with
the usual foliated discrepancy of the foliated MMP.  The theory also gives a
comparison through common adapted refinements: two adapted models with a
common logarithmic simple adapted refinement and the same image separatrix
system define the same adapted toroidal tangential invariants.

The applications include tangential toroidal inversion of adjunction,
branch--conductor descriptions of tangential non-lc and non-klt loci, a
cylinder criterion for tangential log canonicity, and lower semicontinuity of
the adapted toroidal tangential minimal log discrepancy.  The latter is a
pointwise invariant computed as the infimum of the ordinary minimal log
discrepancies of the branch--conductor adjunction pairs lying over the point;
to the author's knowledge, this gives the first arc-space lower
semicontinuity statement for a foliated mld-type invariant in this toroidal
non-resonant sector.

The Mather--Jacobian part is formulated relative to the canonical image
separatrix system associated with the fixed adapted toroidal category.  This
is necessary because non-dicriticality gives a finite formal separatrix
picture, but formal separatrices of simple foliation singularities need not
all be convergent \cite{Spicer}.  Under an algebraic separatrix hypothesis,
the relative Mather--Jacobian invariant agrees with the intrinsic downstairs
separatrix-system invariant.

The hypotheses are sharp in two elementary senses.  Positive resonance among
the logarithmic residues produces reduced tangential arcs whose generic point
is not contained in \(\Dinv\), already in the pure logarithmic model.  Also,
finite-level branch decompositions fail even for the normal crossing divisor
\(xy=0\).  These two phenomena explain the use of non-resonance and reduced
infinite arcs in the main representability theorem.

\subsection*{Main results}

Let $(X,\cF,\Delta)$ be a normal $\bQ$-factorial threefold with a co-rank one
foliated pair.  We work in the logarithmic simple adapted setting described
in Section~\ref{sec:range}.  Fix a logarithmic adapted model
\[
\pi_0:(W_0,\cG_0,\Delta_{W_0})\to (X,\cF,\Delta).
\]
All model independence statements are made inside the toroidal category
$\Adm(X,\cF,\Delta;W_0)$ generated from $W_0$ by coordinate stratum blow-ups.

\begin{Theorem}
\label{thm:intro-canonical-model}
Let $\Sigmatan$ be a tangential centre on an object $W$ of
$\Adm(X,\cF,\Delta;W_0)$.  The reduced tangential sector centred at
$\Sigmatan$ is presented, at the reduced arc level, by the ordinary arc
spaces of the normalised separatrix--conductor system.  More precisely, there is a coequaliser diagram
of reduced pro-functors
\[
\begin{tikzcd}[column sep=large]
\displaystyle\bigsqcup_{\alpha<\beta}J_\infty(C_{\alpha\beta}^{\nu})_{\red}
\arrow[r, shift left=.45ex]
\arrow[r, shift right=.45ex]
&
\displaystyle\bigsqcup_\alpha J_\infty(S_\alpha^\nu)_{\red}
\arrow[r]
&J_\infty(\sS_W^{\sn})_{\red}
\end{tikzcd}
\]
\[
J_\infty(\sS_W^{\sn})_{\red}
\longrightarrow
J_\infty^{\tan}(W,\cG;\Sigmatan)_{\red}.
\]
The first two arrows are induced by the two maps from a normalised conductor
to its adjacent normalised branches.  The presentation identifies adapted
toroidal tangential divisorial cylinders with ordinary maximal divisorial sets
on branch or conductor strata and preserves the logarithmic codimensions
computed from normalised foliated adjunction.  No invariance of ordinary
codimension alone is asserted.
\end{Theorem}

\begin{Theorem}
\label{thm:intro-formal}
Let $(W,\cG)$ be logarithmic simple adapted, let $\Dinv$ be the reduced
invariant separatrix divisor, and set
$\Sigmatan:=\Sing\cG\cap |\Dinv|$.  The reduced tangential arc
pro-functor centred on $\Sigmatan$ is identified with the reduced arc
functor of $\Dinv$ along $\Sigmatan$:
\[
J_\infty^{\tan}(W,\cG;\Sigmatan)_{\red}=J_\infty(\Dinv;\Sigmatan)_{\red}.
\]
Consequently, every irreducible tangential cylinder has a unique generic
branch unless its generic arc lies in a conductor stratum.
\end{Theorem}

\begin{Theorem}
\label{thm:intro-system}
For every object $W$ in the fixed adapted toroidal category, the invariant
separatrices define a finite normalised system consisting of branch pairs $
\{(S_\alpha^\nu,\Theta_\alpha)\}_\alpha
$
and conductor pairs
$
\{(C_{\alpha\beta}^\nu,\Theta_{\alpha\beta})\}_{\alpha<\beta}.$
The system is functorial under coordinate adapted blow-ups in this toroidal
category, and any two objects of the category admit a common adapted toroidal
refinement.  The adjunction boundaries are crepant compatible under these
refinements.
\end{Theorem}

\begin{Corollary}[Tangential Ein--Musta\c{t}\u{a}--Yasuda formula]
\label{thm:intro-emy}
Let $E$ be an adapted toroidal tangential divisorial datum and let
$N_q^{\tan}(E)$ be its divisorial tangential cylinder.  If $E$ is represented
by a divisor $F$ over a normalised branch or conductor stratum $V$, then
\[
\lcodim_{\tan}\bigl(N_q^{\tan}(E)\bigr)=q\,a(F;V,B_V)=q\,a_{\tan}(E;X,\cF,\Delta).
\]
The value is independent of the representative in $\Adm(X,\cF,\Delta;W_0)$
and conductor data are counted once.
\end{Corollary}

\begin{Corollary}[Comparison with foliated discrepancies]
\label{thm:intro-foliated-comparison}
Let $E$ be an adapted toroidal invariant divisor whose tangential datum is
read on a normalised invariant branch.  Then
\[
a_{\tan}(E;X,\cF,\Delta)=\AF(E;X,\cF,\Delta),
\]
where $\AF$ denotes the foliated log discrepancy in the convention of the
foliated MMP.  The comparison is not asserted for conductor data or for
arbitrary divisorial valuations over $X$.
\end{Corollary}

\begin{Theorem}
\label{thm:intro-ioa}
For every closed tangent set $Y\subset X$,
\[
\tmldtor(Y;X,\cF,\Delta)=
\inf_{V\in\mathcal S_{\mathrm{tor}}(X,\cF)}\mld(Y_V;V,B_V),
\]
where the infimum ranges over the adapted toroidal branch and conductor data
of the normalised separatrix system.  Thus the invariant constructed here is
the adapted toroidal tangential mld; toroidal tangential log canonicity,
toroidal tangent lc centres and the toroidal tangential non-klt locus in this
toroidal sector are detected on the normalised branch and conductor adjunction
pairs.
\end{Theorem}

\begin{Theorem}
\label{thm:intro-mj}
For adapted toroidal tangential data whose generic branch descends to a fixed
component of the canonical image separatrix system on $X$, the boundary is
encoded by a formal product of ideals $\bbound_V$, and
\[
a^{\tan}_{\MJ}(E;X,\cF,\Delta)=a_{\MJ}(F;V,\bbound_V)
\]
is independent of the resolved branch presentation.  The corresponding
Mather--Jacobian logarithmic codimension formula is
\[
\lcodim^{\MJ}_{\tan}\bigl(N_q^{\tan}(E)\bigr)=q\,a^{\tan}_{\MJ}(E;X,\cF,\Delta).
\]
In the locally complete intersection range with trivial Jacobian correction
this reduces to the ordinary tangential theory.
\end{Theorem}

\subsection*{Structure of the paper}

Sections~\ref{sec:range}--\ref{sec:formal} fix the logarithmic simple
adapted setting and prove the reduced arc-level tangential confinement
statement.  This part also records the sharpness of the non-resonance
condition and the finite-level obstruction to a global branch decomposition.
Sections~\ref{sec:category}--\ref{sec:adjunction} construct the normalised
branch--conductor system in the fixed adapted toroidal category, prove its
functoriality, and establish normalised foliated adjunction with conductor
compatibility.  The comparison through common adapted refinements is also
proved there.
Sections~\ref{sec:valuations}--\ref{sec:foliated-comparison} define adapted
toroidal tangential divisorial data and compare the tangential discrepancy
with the usual foliated discrepancy of the foliated MMP for adapted invariant
branch data. 
Section~\ref{sec:emy} proves the tangential Ein--Musta\c{t}\u{a}--Yasuda
formula.  Sections~\ref{sec:ioa}--\ref{sec:transverse} apply this formula to
tangential inversion of adjunction, non-lc and non-klt loci, cylinder
criteria, lower semicontinuity, and the adapted tangent/transverse
decomposition. 
Section~\ref{sec:mj} gives the Mather--Jacobian refinement on the canonical
image separatrix system on $X$ and records the intrinsic comparison under the
algebraic separatrix hypothesis.


\section{The logarithmic simple adapted setting}
\label{sec:range}

 Throughout the article, we work over the field \(\bC\).  A foliation means a. saturated involutive subsheaf \(\cF\subset T_X\).  A co-rank one foliation on a smooth threefold is locally given by an integrable \(1\)-form, uniquely
determined up to multiplication by a unit.

\begin{Definition}[Logarithmic simple adapted chart]
\label{def:simple-chart}
Let \(W\) be a smooth threefold and let \(\mathcal G\) be a co-rank one
foliation.  A point \(p\in W\) has a \emph{logarithmic simple adapted chart}
if there are formal coordinates \(x_1,x_2,x_3\) centred at \(p\), a reduced
invariant normal crossing divisor
\[
\Dinv=\{x_1\cdots x_r=0\},
\qquad 1\le r\le 3,
\]
and a local generator
\[
\omega=
u\,x_1\cdots x_r
\left(
\sum_{i=1}^r\lambda_i\frac{dx_i}{x_i}
+
\sum_{j=r+1}^3 h_j(x)\,dx_j
\right),
\]
where \(u\) is a unit and \(\lambda_i\ne0\), such that
\[
a_1\lambda_1+\cdots+a_r\lambda_r\ne0
\]
for every non-zero vector \((a_1,\ldots,a_r)\in\bZ_{\geq 0}^r\).

The adapted divisor \(\Dinv\) is required to contain the invariant
separatrix branches relevant to the tangential centre under consideration.
For the intrinsic separatrix sector, the additional local separatrix hypothesis is imposed:
these branches are required to be precisely the formal
separatrices through the centre.  The remaining boundary components of
\(\Delta_W\) passing through \(p\) are required to be among the coordinate
hyperplanes, and \(\Dinv+\operatorname{Supp}\Delta_W\) is required to
have simple normal crossings.
\end{Definition}



\begin{Definition}
An \emph{adapted blow-up} of a simple adapted model is the blow-up of a smooth stratum of $\Dinv+\Supp\Delta_W$ whose generic point is either tangent to $\cG$ or transverse to it.  In the tangential part of the article, the construction uses the tangent adapted blow-ups, namely those whose centre is contained in $\Dinv$.
\end{Definition}

\begin{Definition}
A foliated pair $(X,\cF,\Delta)$ is in the \emph{logarithmic simple adapted setting} if it admits a projective birational morphism
\[
\pi:(W,\cG,\Delta_W)\to (X,\cF,\Delta)
\]
with $W$ smooth such that every point of the tangential locus of interest on $W$ has a logarithmic simple adapted chart and $\Dinv+\Supp\Delta_W$ is SNC.  All global statements in what follows are made within this setting.  Thus the setting is a chosen adapted part of the simple models produced after reduction of singularities; it is not asserted that every canonical or log canonical foliation is already of this form.
\end{Definition}


\begin{Remark}[Relation with foliated MMP singularities]
The logarithmic simple adapted setting used in this article is modelled on the
simple singularities of codimension-one foliations appearing in the foliated
minimal model programme.  Cano's reduction theorem states that, on a smooth
threefold, a codimension-one foliation admits a resolution by blow-ups with
centres in the singular locus of the foliation such that the transformed
foliation has only simple singularities \cite{Can04}.  Thus the passage to
simple normal forms is resolution-theoretic in origin: it is obtained on a
birational model by explicit blow-ups, not by replacing the foliation with an
unrelated rational normal form.
The full class of simple singularities is not used.  In the usual
definition there is a logarithmic type
\[
\omega=(x_1\cdots x_r)\left(\sum_{i=1}^r
\lambda_i\frac{dx_i}{x_i}\right),
\]
with the condition that
\[
\sum_{i=1}^r a_i\lambda_i=0
\qquad (a_i\in\bZ_{\geq 0})
\]
forces all \(a_i\) to vanish, and there is also a second type involving a
non-unit function \(\psi\) \cite[Definition~2.13]{Spicer}.  We work in the logarithmic non-resonant part.  This is exactly the part in
which the residue computation used in what follows confines reduced tangential arcs to
the adapted invariant divisor \(\Dinv\).  The second simple type is not
covered by the arc-confinement argument in this article.

Simple singularities are canonical, and hence log canonical, but the converse
is false \cite[Lemma~2.15 and Example~2.16]{Spicer}.  Consequently one cannot
replace the non-resonant logarithmic normal form by the assumption that the
foliated pair is canonical or log canonical.  The role of positive
non-resonance is different: it excludes resonant tangential arcs
transverse to the chosen invariant divisor and is precisely what allows the
reduced tangential arc sector to be confined to \(\Dinv\).  This
confinement is what permits the subsequent discrepancy calculation to be
reduced to ordinary arc-space discrepancy formulas on the branch--conductor
adjunction pairs.

Non-dicriticality enters separately.  It guarantees that exceptional divisors
over tangent centres remain invariant and, in the simple setting, gives a
finite separatrix picture \cite[Definition~2.17--Example~2.22]{Spicer}.
However, even for simple foliation singularities, formal separatrices need
not all be convergent \cite[p.~8]{Spicer}.  This is why the
Mather--Jacobian refinement in what follows is formulated relative to the canonical
image separatrix system associated with the fixed adapted toroidal category,
rather than as an unconditional statement about all intrinsic formal
separatrices downstairs.  Cano--Cerveau's separatrix results explain the
analytic relevance of non-dicriticality, but they do not remove the need for
an algebraic or image-system formulation in the Mather--Jacobian part
\cite{CC92}.
\end{Remark}

\section{The adapted model category and the basic objects}
\label{sec:category}



\begin{Definition}[The adapted toroidal category]
Fix a logarithmic adapted model $W_0\to X$. Let $\Adm(X,\cF,\Delta;W_0)$ be the toroidal category whose objects are obtained from $W_0$ by finite sequences of coordinate stratum blow-ups, and write $\Adm(X,\cF,\Delta)$ when $W_0$ is understood. Its objects are adapted toroidal models
\[
(W,\cG,\Dinv,\Delta_W)\to (X,\cF,\Delta)
\]
in the logarithmic simple adapted setting.  A morphism
\[
\mu:(W',\cG',\Dinv',\Delta_{W'})\longrightarrow (W,\cG,\Dinv,\Delta_W)
\]
is a composition of blow-ups of smooth strata of $\Dinv+\Supp\Delta_W$ such that
\[
\Dinv'=(\mu^{-1}\Dinv)_{\red},
\qquad
K_{\cG'}+\Delta_{W'}=\mu^*(K_{\cG}+\Delta_W)+\sum_E \alpha(E;\cG,\Delta_W)E,
\]
and each intermediate model remains in the logarithmic simple adapted setting.  The positive non-resonance condition in Definition~\ref{def:simple-chart} is part of the data and guarantees that residue sums along exceptional invariant components do not vanish.
The full adapted toroidal category allows tangent and transverse coordinate
stratum blow-ups.  The tangential discrepancy theory below uses the tangent
part of this category; transverse strata enter through boundary and comparison
data.
\end{Definition}


\begin{Notation}[Discrepancy conventions]
In the relative pull-back formula above, \(\alpha(E;\cG,\Delta_W)\) denotes
the coefficient of \(E\) in the relative foliated canonical divisor.  Foliated
log discrepancies are denoted by \(\AF\).  In the local tangent toroidal
calculations below, the quantities written \(\AF(E;W,\cG,\Delta_W)\) are in
this log-discrepancy normalisation and are compared directly with the ordinary
log discrepancies of the branch adjunction pairs.  Thus relative coefficients
and log discrepancies are not interchanged.
\end{Notation}


\begin{Definition}
For $W\in\Adm(X,\cF,\Delta)$, let $\Strata(W)$ be the finite set consisting of:
\begin{enumerate}[label=\textup{(\roman*)},leftmargin=2.4em]
\item the normalisations $S_\alpha^\nu$ of the irreducible invariant branches $S_\alpha\subset \Dinv$;
\item the normalisations $C_{\alpha\beta}^\nu$ of the reduced pairwise conductor curves $S_\alpha\cap S_\beta$.
\end{enumerate}
Triple points are not extra strata; they are closed points of the normalised conductor curves and are resolved by adapted toroidal refinements.
\end{Definition}

\begin{Remark}
On the adapted smooth models, the local SNC branches are already normal.  The
notation \(S_\alpha^\nu\) is retained in order to keep track of the
functorial normalised branch--conductor system and of its image downstairs.
\end{Remark}

\begin{Construction}[Seminormal pushout and equaliser description]
\label{constr:sn-pushout-equaliser}
For \(W\in\Adm(X,\cF,\Delta)\), let \(\sS_W^{\sn}\) be the seminormal
pushout of the branch--conductor diagram
\[
\begin{tikzcd}[column sep=large]
\displaystyle\bigsqcup_{\alpha<\beta} C_{\alpha\beta}^{\nu}
\arrow[r, shift left=.45ex]
\arrow[r, shift right=.45ex]
&
\displaystyle\bigsqcup_{\alpha} S_\alpha^\nu
\arrow[r]
&
\sS_W^{\sn}.
\end{tikzcd}
\]
The two arrows are induced by the maps from a normalised conductor to its
two adjacent normalised branches,
\[
C_{\alpha\beta}^\nu\longrightarrow S_\alpha^\nu,
\qquad
C_{\alpha\beta}^\nu\longrightarrow S_\beta^\nu .
\]

Equivalently, the structure sheaf of \(\sS_W^{\sn}\) is described by the
equaliser
\begin{equation}\label{eq:sn-equaliser-operational}
\cO_{\sS_W^{\sn}}
=
\ker\left(
\begin{tikzcd}[ampersand replacement=\&, column sep=large, baseline=(current bounding box.center)]
\displaystyle\bigoplus_\alpha (\iota_\alpha)_*\cO_{S_\alpha^\nu}
\arrow[r, shift left=.45ex]
\arrow[r, shift right=.45ex]
\&
\displaystyle\bigoplus_{\alpha<\beta}(\iota_{\alpha\beta})_*\cO_{C_{\alpha\beta}^\nu}
\end{tikzcd}
\right).
\end{equation}
Thus sections of \(\cO_{\sS_W^{\sn}}\) are precisely collections
\((s_\alpha)_\alpha\) with
$
s_\alpha|_{C_{\alpha\beta}^\nu}
=
s_\beta|_{C_{\alpha\beta}^\nu}$
for all $\alpha<\beta$.

\end{Construction}

\begin{Lemma}[Existence and reducedness of the pushout]
\label{lem:sn-pushout}
In the logarithmic SNC charts used in this article, \eqref{eq:sn-equaliser-operational} defines a reduced seminormal scheme finite over $\Dinv\subset W$.  It is the pushout of the finite branch--conductor diagram in the category of seminormal schemes.
\end{Lemma}

\begin{proof}
The assertion is local on $W$.  At a double point the model is
\[
D=(xy=0)\subset \Spec A,
\qquad
S_x=\Spec A/(x),\quad S_y=\Spec A/(y),\quad C=\Spec A/(x,y).
\]
The fibre product description of the coordinate ring is
\[
A/(xy)\simeq A/(x)\times_{A/(x,y)}A/(y),
\]
because an element of the right-hand side is a pair $(\bar f,\bar g)$ with $\bar f|_C=\bar g|_C$, and it is represented by $g+xh$ after lifting $\bar f-\bar g$ to $xh$.  This is exactly the equaliser \eqref{eq:sn-equaliser-operational}.  At a triple point the model is $D=(xyz=0)$.  The seminormal gluing is the equaliser imposing all three pairwise conductor compatibilities:
\[
A/(xyz)\simeq
\left\{(f_x,f_y,f_z)\in A/(x)\oplus A/(y)\oplus A/(z)\;\middle|\;
\begin{array}{l}
f_x|_{(x,y)}=f_y|_{(x,y)},\\
f_x|_{(x,z)}=f_z|_{(x,z)},\\
f_y|_{(y,z)}=f_z|_{(y,z)}
\end{array}
\right\}.
\]
Indeed, the map $A/(xyz)$ to the right-hand side is injective because the components cover $D$, and surjectivity is the elementary Chinese-remainder calculation for the square-free monomial ideal $(xyz)$ with its three conductor ideals.  Since the components and conductor curves are reduced and the gluing is along reduced finite closed subschemes, the resulting pushout is reduced and seminormal.  The general SNC case is obtained by adjoining smooth parameters.
\end{proof}

\begin{Proposition}
\label{prop:separatrix-functor}
The construction
$
W\longmapsto \sS_W^{\sn}
$
defines a contravariant functor
\[
\sS:\Adm(X,\cF,\Delta)^{\mathrm op}\longrightarrow \SN,
\]
where $\SN$ is the category of reduced seminormal schemes finite over the relevant tangential centre.  For a morphism $\mu:W'\to W$, the map $\sS_{W'}^{\sn}\to\sS_W^{\sn}$ is induced on normalisations by strict transform maps and, for exceptional branches over conductor curves, by the projection to the old normalised conductor.
\end{Proposition}

\begin{proof}
It suffices to treat the case in which
$
\mu:(W',\mathcal G',\Delta_{W'})\longrightarrow
(W,\mathcal G,\Delta_W)
$
is one adapted coordinate blow-up. The general case then follows by composition.
Write the normalised branch--conductor diagrams of \(W\) and \(W'\) as
\[
\mathfrak D_W=
\left(
\bigsqcup_{\alpha<\beta} C_{\alpha\beta}^{\nu}
\rightrightarrows
\bigsqcup_{\alpha} S_\alpha^{\nu}
\right),
\qquad
\mathfrak D_{W'}=
\left(
\bigsqcup_{\gamma<\delta} C_{\gamma\delta}^{\prime\,\nu}
\rightrightarrows
\bigsqcup_{\gamma} S_\gamma^{\prime\,\nu}
\right).
\]
It remains to construct a morphism of diagrams
$
\mathfrak D_{W'}\longrightarrow \mathfrak D_W.
$
The induced morphism on the seminormal pushouts then follows from the
equaliser description
\[
\mathcal O_{\mathscr S_W^{\sn}}
=
\ker\left(
\bigoplus_\alpha (\iota_\alpha)_*\mathcal O_{S_\alpha^\nu}
\rightrightarrows
\bigoplus_{\alpha<\beta}(\iota_{\alpha\beta})_*
\mathcal O_{C_{\alpha\beta}^\nu}
\right).
\tag{\ref{eq:sn-equaliser-operational}}
\]
First consider a branch \(S_\gamma'\subset D'_{\inv}\). There are three possibilities.

If \(S_\gamma'\) is the strict transform of a branch \(S_\alpha\), then
\(\mu(S_\gamma')\subset S_\alpha\), and by the universal property of
normalisation one obtains a unique morphism
$
\widetilde\mu_\gamma:S_\gamma^{\prime\,\nu}\longrightarrow S_\alpha^\nu
$
making the diagram commute:
\[
\begin{tikzcd}
S_\gamma^{\prime\,\nu} \arrow[r] \arrow[d,"\widetilde\mu_\gamma"']&
S_\gamma' \arrow[d,"\mu|_{S_\gamma'}"]\\
S_\alpha^\nu \arrow[r]& S_\alpha.
\end{tikzcd}
\]

If the centre of \(\mu\) is a conductor curve
$
C_{\alpha\beta}=S_\alpha\cap S_\beta,
$
then the exceptional branch \(E\subset W'\) is a ruled surface over the
normalisation of the conductor:
\[
\pi_E:E^\nu\longrightarrow C_{\alpha\beta}^{\nu}.
\]
Thus the map attached to the exceptional branch is
\[
E^\nu \xrightarrow{\ \pi_E\ } C_{\alpha\beta}^{\nu}
\longrightarrow \mathscr S_W^{\sn}.
\]
In particular, \(E^\nu\) is not identified birationally with
\(C_{\alpha\beta}^{\nu}\); rather, it is a new branch of the refined
system whose structural map factors through the old conductor stratum.
If the centre is a triple point
$
p=S_\alpha\cap S_\beta\cap S_\eta,
$
then, in local coordinates \(\Dinv=(xyz=0)\), the blow-up is the star
subdivision of the cone
\[
\sigma=\mathbb R_{\geq 0}\langle e_x,e_y,e_z\rangle
\]
by the ray \(e_x+e_y+e_z\). The exceptional branch \(E\) maps to \(p\).
Each new conductor curve corresponds to an edge of the subdivided incidence
complex, and its image is the old stratum determined by the minimal face of
\(\sigma\) containing that edge. Thus a new conductor maps either to one of
\[
C_{\alpha\beta}^{\nu},\quad
C_{\alpha\eta}^{\nu},\quad
C_{\beta\eta}^{\nu},
\]
or to the point \(p\), according to its incidence.
It remains to check compatibility on conductor strata. Let
$
C_{\gamma\delta}'=S_\gamma'\cap S_\delta'
$ 
be a conductor curve of \(D'_{\inv}\). Suppose that the maps constructed
above send \(S_\gamma^{\prime\,\nu}\) and \(S_\delta^{\prime\,\nu}\) to
strata \(V_\gamma,V_\delta\) of \(\mathfrak D_W\). Since
$
\mu(C_{\gamma\delta}')\subset
\mu(S_\gamma')\cap\mu(S_\delta'),
$
the normalisation \(C_{\gamma\delta}^{\prime\,\nu}\) maps, by the universal
property of normalisation, to the normalised conductor stratum of
\(\mathfrak D_W\) through which both restrictions factor. Equivalently, the
following diagram commutes:
\[
\begin{tikzcd}
C_{\gamma\delta}^{\prime\,\nu}
\arrow[r] \arrow[d] &
S_\gamma^{\prime\,\nu}
\arrow[d]\\
C_{\alpha\beta}^{\nu}
\arrow[r] &
S_\alpha^\nu
\end{tikzcd}
\qquad
\begin{tikzcd}
C_{\gamma\delta}^{\prime\,\nu}
\arrow[r] \arrow[d] &
S_\delta^{\prime\,\nu}
\arrow[d]\\
C_{\alpha\beta}^{\nu}
\arrow[r] &
S_\beta^\nu.
\end{tikzcd}
\]
In this case \(C_{\alpha\beta}^{\nu}\) may be replaced by a point stratum in the
triple-point case. Hence the two conductor restriction maps agree after
composition with the corresponding maps in \(\mathfrak D_W\).

Consequently, the morphisms on the normalised branches and conductors define a
morphism of diagrams
\[
\begin{tikzcd}
\displaystyle\bigsqcup_{\gamma<\delta} C_{\gamma\delta}^{\prime\,\nu}
\arrow[r,shift left=.6ex] \arrow[r,shift right=.6ex] \arrow[d] &
\displaystyle\bigsqcup_{\gamma} S_\gamma^{\prime\,\nu}
\arrow[d]\\
\displaystyle\bigsqcup_{\alpha<\beta} C_{\alpha\beta}^{\nu}
\arrow[r,shift left=.6ex] \arrow[r,shift right=.6ex] &
\displaystyle\bigsqcup_{\alpha} S_\alpha^\nu.
\end{tikzcd}
\]
Passing to structure sheaves gives a commutative diagram
\[
\begin{tikzcd}
\displaystyle
\bigoplus_\alpha(\iota_\alpha)_*\mathcal O_{S_\alpha^\nu}
\arrow[r,shift left=.6ex] \arrow[r,shift right=.6ex] \arrow[d] &
\displaystyle
\bigoplus_{\alpha<\beta}
(\iota_{\alpha\beta})_*\mathcal O_{C_{\alpha\beta}^\nu}
\arrow[d]\\
\displaystyle
\bigoplus_\gamma(\iota'_\gamma)_*
\mathcal O_{S_\gamma^{\prime\,\nu}}
\arrow[r,shift left=.6ex] \arrow[r,shift right=.6ex] &
\displaystyle
\bigoplus_{\gamma<\delta}
(\iota'_{\gamma\delta})_*
\mathcal O_{C_{\gamma\delta}^{\prime\,\nu}}.
\end{tikzcd}
\]
Thus a section of
\(\mathcal O_{\mathscr S_W^{\sn}}\), viewed as a compatible tuple of
sections on the normalised branches \(S_\alpha^\nu\), pulls back to a
compatible tuple on the refined branches \(S_\gamma^{\prime\,\nu}\).
By the equaliser formula, this defines a morphism
\[
\mathscr S_{W'}^{\sn}\longrightarrow \mathscr S_W^{\sn}.
\]
Finally, suppose that
\[
W''\xrightarrow{\mu'} W'\xrightarrow{\mu} W
\]
are two adapted coordinate blow-ups, the morphism attached to
\(\mu\circ\mu'\) is the same as the composite of the morphisms attached to
\(\mu'\) and \(\mu\). Indeed, on every normalised branch and conductor this
is the uniqueness part of the universal property of normalisation, and on
the seminormal pushout it follows from the universal property of the
equaliser. Hence the construction is functorial.
\end{proof}

\begin{Definition}
\label{def:tangential-cylinder-operational}
A \emph{tangential divisorial cylinder} is an equivalence class of triples
$
(W,V,C_V),
$
where $W\in\Adm(X,\cF,\Delta)$, $V\in\Strata(W)$, and $C_V\subset J_\infty(V)$ is an ordinary maximal divisorial cylinder.  Two triples $(W,V,C_V)$ and $(W',V',C_{V'})$ are equivalent if on some common adapted toroidal refinement $U$ their pull-backs to the corresponding strict or exceptional strata of $\sS_U^{\sn}$ coincide.
The ordinary codimension $\codim_{J_\infty(V)}(C_V)$ is attached to the representative.  The invariant quantity is the logarithmic codimension after the adjunction boundary $B_V$ has been fixed:
\[
\lcodim_{\tan}(W,V,C_V):=\codim_{J_\infty(V)}(C_V)-\ord_{C_V}(B_V).
\]
\end{Definition}

\begin{Lemma}[Independence of tangential logarithmic codimension]
\label{lem:codim-independence}
If $(W,V,C_V)$ and $(W',V',C_{V'})$ are equivalent adapted toroidal divisorial cylinders, then
\[
\lcodim_{\tan}(W,V,C_V)=\lcodim_{\tan}(W',V',C_{V'}).
\]
\end{Lemma}
\begin{proof}
Choose a common adapted toroidal refinement and let $U_V$ be the normalised stratum carrying the common maximal divisorial set $C_U$.  Write \(p:U_V\to V\) and \(q:U_V\to V'\) for the induced birational maps.  Crepant compatibility of normalised foliated adjunction gives
\[
K_{U_V}+B_{U_V}=p^*(K_V+B_V)+A_p=q^*(K_{V'}+B_{V'})+A_q,
\]
where $A_p$ and $A_q$ have the same coefficient on the divisor defining $C_U$.
Let $C_U=N_q(F)$ be the maximal divisorial set associated with a prime divisor
$F$ on $U_V$ and multiplicity $q$, namely the usual divisorial cylinder of arcs
with order $q$ along $F$. Since $V$ and $V'$ are smooth at the generic point
relevant to the adapted toroidal datum, the ordinary maximal-divisorial-set
formula gives
\[
\lcodim_{(V,B_V)}(C_V)=q\bigl(k_F(V)+1-\ord_F(B_V)\bigr),
\]
and similarly over $V'$.  Pulling the crepant equality to $U_V$ gives
\[
k_F(V)+1-\ord_F(B_V)=k_F(V')+1-\ord_F(B_{V'}).
\]
It follows that the logarithmic codimension is invariant.  No invariance of the ordinary codimension alone is asserted.
\end{proof}

\subsection{Comparison through common adapted refinements}
\label{subsec:common-refinement}

The construction above is made inside an adapted toroidal category generated
from an initial model.  The following result records the conditional
independence which is used when two initial choices have a common adapted
refinement and determine the same image separatrix system on $X$.

\begin{Theorem} 
\label{thm:common-refinement-invariance}
Let $W_1\to (X,\cF,\Delta)$ and $W_2\to (X,\cF,\Delta)$ be logarithmic
simple adapted non-dicritical models.  Assume that there is a logarithmic
simple adapted model $U$ with adapted toroidal morphisms $U\to W_1$ and
$U\to W_2$.  Assume moreover that $W_1$ and $W_2$ induce the same canonical
image separatrix system on $X$.  Then their adapted toroidal tangential
divisorial data are identified after pull-back to $U$, and their tangential
discrepancies agree.  Consequently, for every closed tangential set $Y$
contained in the common image system,
\[
\tmld_{\mathrm{tor}}^{W_1}(Y;X,\cF,\Delta)
=
\tmld_{\mathrm{tor}}^{W_2}(Y;X,\cF,\Delta).
\]
\end{Theorem}

\begin{proof}
Let $\rho_i:U\to W_i$ be the two adapted toroidal morphisms.  By
functoriality of the seminormal branch--conductor construction, each
$\rho_i$ induces a morphism
\[
\sS_U^{\sn}\longrightarrow \sS_{W_i}^{\sn}.
\]
Thus every adapted toroidal tangential divisorial datum represented on
$W_i$ pulls back to a divisorial datum on a normalised branch or conductor
stratum of $\sS_U^{\sn}$.

Let a datum on $W_1$ be represented by $N_q(F_1)\subset J_\infty(V_1)$,
where $V_1$ is a normalised branch or conductor stratum of
$\sS_{W_1}^{\sn}$.  Its pull-back to $U$ is represented by
$N_q(F_U)\subset J_\infty(V_U)$ for a corresponding normalised stratum
$V_U$ of $\sS_U^{\sn}$.  By crepant compatibility of the normalised
adjunction boundaries and Lemma~\ref{lem:codim-independence},
\[
a(F_1;V_1,B_{V_1})=a(F_U;V_U,B_{V_U}).
\]
The same construction applies to $W_2$.  Since $W_1$ and $W_2$ induce the
same canonical image separatrix system on $X$, the data pulled back to $U$
represent the same image branch or conductor datum downstairs.  Thus the
corresponding datum on $W_2$ has the same pull-back $F_U$, and hence
\[
a(F_1;V_1,B_{V_1})=a(F_U;V_U,B_{V_U})=a(F_2;V_2,B_{V_2}).
\]
Equivalently, $a_{\tan}^{W_1}=a_{\tan}^{W_2}$ on corresponding adapted
toroidal tangential data.  Taking the infimum over all such data with centre
contained in $Y$ gives the asserted equality of $\tmld_{\mathrm{tor}}$.
\end{proof}

 

\section{Formal tangential representability}
\label{sec:formal}

Carter's jet schemes of a foliation \cite{Carter,CarterThesis} provide the
finite-jet language motivating the tangential sector.  The discrepancy
calculations in what follows, however, use only the reduced infinite tangential arc
functor.  Thus, when a co-rank one foliation is locally generated by an
integrable one-form $\omega$, the reduced functor is used
\[
J_\infty^{\tan}(W,\cG;Z)_{\red}
=
\{\gamma\in J_\infty(W)_{\red}:\gamma(0)\in Z,\ \gamma^*\omega=0\}.
\]
This condition is independent of the chosen local generator, since
multiplying $\omega$ by a unit does not change the vanishing of its pull-back.
This is the precise arc-level consequence of Carter's tangent-jet viewpoint
which is needed in the sequel.  The finite tangent jet schemes may contain
nilpotent and truncation dependent structure; no equality of finite-level
Carter jet schemes is used in any cylinder or codimension calculation in this
article.

The next theorem is the only reduced Carter-type result required in what follows.  It is
best regarded as an arc-confinement statement rather than as a finite-jet
representability theorem.

\begin{Theorem} 
\label{thm:carter-representability}
Let $p\in W$ have a logarithmic simple adapted chart and write
\[
R=\bC[[x_1,x_2,x_3]],\qquad f=x_1\cdots x_r,
\qquad \Dinv=(f=0).
\]
Let
\[
\omega=
u\,x_1\cdots x_r
\left(
\sum_{i=1}^r \lambda_i\frac{dx_i}{x_i}
+
\sum_{j=r+1}^3 h_j(x)dx_j
\right)
\]
be the local generator, where $u$ is a unit and the residues satisfy the
positive non-resonance condition.  Then, for every closed subset
$Z\subset \Dinv$, the reduced tangential arc functor centred on $Z$
coincides with the reduced arc functor of $\Dinv$ centred on $Z$.
Equivalently, every reduced tangential formal arc centred on $Z$ factors
through the adapted invariant divisor $\Dinv$, and every arc on
$\Dinv$ is tangential.  In particular, for
$\Sigmatan=\Sing\cG\cap |\Dinv|$,
\[
J_\infty^{\tan}(W,\cG;\Sigmatan)_{\red}
=
J_\infty(\Dinv;\Sigmatan)_{\red}.
\]
This is an equality of reduced arc functors only.  No assertion is made about
the non-reduced finite Carter jet schemes, nor about equality of finite-level
reduced $m$-jet schemes for all $m$.
\end{Theorem}

 
\begin{proof}
The assertion is local at \(p\).  Since multiplication by a unit does not
change the vanishing of the pull-back of a one-form, we may omit \(u\).
Let
\[
\gamma:\Spec K[[t]]\longrightarrow W,
\qquad
\gamma(t)=(x_1(t),x_2(t),x_3(t)),
\]
be a formal arc centred on \(\Dinv\).  We first prove that a tangential arc
centred on \(\Dinv\) is contained in \(\Dinv\).

Assume, by contradiction, that \(\gamma\) is not contained in \(\Dinv\).
Then none of \(x_1(t),\ldots,x_r(t)\) is the zero series.  Write
\[
a_i=\ord_t x_i(t),\qquad
x_i(t)=t^{a_i}u_i(t),
\qquad
u_i(0)=c_i\in K^*,
\]
and set
\[
P(t)=x_1(t)\cdots x_r(t),\qquad
A=\ord_t P(t)=\sum_{i=1}^r a_i .
\]
Since the closed point of \(\gamma\) lies on \(\Dinv\), at least one \(a_i\)
is positive.  Thus \(a=(a_1,\ldots,a_r)\) is a non-zero element of
\(\bZ_{\geq 0}^r\).
Pulling back the logarithmic generator gives
\[
\gamma^*\omega
=
P(t)
\left(
\sum_{i=1}^r\lambda_i\frac{x_i'(t)}{x_i(t)}
+
\sum_{j=r+1}^3 h_j(\gamma(t))x_j'(t)
\right)dt .
\]
We analyse the two summands separately.  Since
\(x_i(t)=t^{a_i}u_i(t)\) with \(u_i(0)\neq0\), one has
\[
\frac{x_i'(t)}{x_i(t)}
=
\frac{a_i}{t}+\frac{u_i'(t)}{u_i(t)}.
\]
Hence
\[
\sum_{i=1}^r\lambda_i\frac{x_i'(t)}{x_i(t)}
=
\left(\sum_{i=1}^r a_i\lambda_i\right)\frac{1}{t}
+
\sum_{i=1}^r\lambda_i\frac{u_i'(t)}{u_i(t)}.
\]
The second term on the right-hand side is regular in \(t\).  Moreover
\[
P(t)
=
t^A\prod_{i=1}^r u_i(t)
=
t^A
\left(
\prod_{i=1}^r c_i+O(t)
\right).
\]
Therefore the logarithmic part contributes
\[
P(t)
\sum_{i=1}^r\lambda_i\frac{x_i'(t)}{x_i(t)}
=
\left(\sum_{i=1}^r a_i\lambda_i\right)
\left(\prod_{i=1}^r c_i\right)t^{A-1}
+
O(t^A).
\]
On the other hand, the regular part satisfies
\[
h_j(\gamma(t))x_j'(t)\in K[[t]]
\]
for every \(j>r\).  Hence
\[
P(t)\sum_{j=r+1}^3 h_j(\gamma(t))x_j'(t)
\in t^A K[[t]].
\]
Thus the coefficient of \(t^{A-1}dt\) in \(\gamma^*\omega\) is exactly
\[
\left(\sum_{i=1}^r a_i\lambda_i\right)
\prod_{i=1}^r c_i .
\]
By positive non-resonance,
\[
\sum_{i=1}^r a_i\lambda_i\neq0,
\]
and since each \(c_i\neq0\), this coefficient is non-zero.  Consequently
\(\gamma^*\omega\neq0\), contradicting tangency.  Hence every reduced
tangential arc centred on \(\Dinv\) is contained in \(\Dinv\).

Conversely, assume that \(\gamma\) factors through \(\Dinv\).  Since
$
\Dinv=(x_1\cdots x_r=0)
$
and \(K[[t]]\) is an integral domain, there exists an index
\(i_0\in\{1,\ldots,r\}\) such that
$
x_{i_0}(t)\equiv0.
$
Write \(f=x_1\cdots x_r\).  The logarithmic form is the regular one-form
\[
\omega
=
u\left(
\sum_{i=1}^r\lambda_i\frac{f}{x_i}\,dx_i
+
f\sum_{j=r+1}^3 h_j\,dx_j
\right).
\]
Pulling back, the second summand vanishes because $
f(\gamma(t))=0.
$
For the first summand, if \(i\neq i_0\), then \(f/x_i\) still contains the
factor \(x_{i_0}\), so
\[
\left(\frac{f}{x_i}\right)(\gamma(t))=0.
\]
If \(i=i_0\), then \(dx_{i_0}(t)=0\), because \(x_{i_0}(t)\equiv0\).
Therefore every term in \(\gamma^*\omega\) vanishes, and hence
$
\gamma^*\omega=0.
$
Thus every arc on \(\Dinv\) is tangential.

We have proved, after arbitrary extension of the ground field, that the
\(K[[t]]\)-points of the reduced tangential arc functor centred on
\(\Dinv\) coincide with the \(K[[t]]\)-points of \(J_\infty(\Dinv)\).  The
same argument applies after imposing the closed condition \(\gamma(0)\in Z\)
for any closed subset \(Z\subset\Dinv\).  Therefore
\[
J_\infty^{\tan}(W,\cG;Z)_{\red}
=
J_\infty(\Dinv;Z)_{\red}.
\]
 
\end{proof}

\subsection{Sharpness of the non-resonance condition}

The positive non-resonance condition is sharp already for pure logarithmic
normal forms.

\begin{Definition} 
For residues $\lambda_1,\ldots,\lambda_r$, set
\[
\Rlambda=
\{a\in\bZ_{\geq 0}^r\setminus\{0\}:a\cdot\lambda=0\}.
\]
\end{Definition}

\begin{Proposition} 
\label{prop:resonant-arcs-pure-log}
Let
\[
\omega=x_1\cdots x_r\sum_{i=1}^r
\lambda_i\frac{dx_i}{x_i}
\]
and let $\Dinv=(x_1\cdots x_r=0)$.  If $\Rlambda\neq\varnothing$, then
the reduced tangential arc sector centred on $\Dinv$ contains arcs whose
generic point is not contained in $\Dinv$.
\end{Proposition}

\begin{proof}
Take $a\in\Rlambda$ and $c_i\in\bC^*$, and set $x_i(t)=c_it^{a_i}$.
Then
\[
\frac{dx_i(t)}{x_i(t)}=a_i\frac{\dd t}{t},
\]
hence
\[
\gamma^*
\left(
\sum_i\lambda_i\frac{dx_i}{x_i}
\right)
=
(a\cdot\lambda)\frac{\dd t}{t}=0.
\]
Thus $\gamma^*\omega=0$.  Since no $x_i(t)$ is identically zero, the
generic point of $\gamma$ is not contained in $\Dinv$.  The closed point
lies on $\Dinv$ because $a\neq0$.
\end{proof}

\begin{Remark}
For the full adapted logarithmic form with regular terms, the condition
$a\in\Rlambda$ is the vanishing of the leading obstruction used in the
confinement theorem.  A complete description of the resonant sector in the
presence of the regular terms requires additional initial-form analysis.
Thus the result above is used only as a sharpness statement, not as a
classification of all resonant tangential arcs.
\end{Remark}

\subsection{A finite-level obstruction}

The reduced infinite-arc confinement theorem should not be replaced by a
global finite-level statement.  The obstruction already appears for an
ordinary normal crossing divisor.

\begin{Example}[Mixed finite jets]
\label{ex:mixed-finite-jets}
Let $D=(xy=0)\subset\bA^2$.  For every $m\geq 1$, the $m$-jet
$x(t)=t$, $y(t)=t^m$ defines a point of $J_m(D)$, but it does not factor
through $J_m(x=0)\cup J_m(y=0)$.  Consequently, no global finite-level branch
decomposition $J_m(D)=J_m(x=0)\cup J_m(y=0)$ can hold.
\end{Example}

\begin{proof}
An $m$-jet is computed modulo $t^{m+1}$.  For the displayed jet,
$x(t)y(t)=t^{m+1}$, hence $x(t)y(t)\equiv0\pmod{t^{m+1}}$.  Thus it
defines a point of $J_m(D)$.  However, $x(t)=t\not\equiv0\pmod{t^{m+1}}$
and $y(t)=t^m\not\equiv0\pmod{t^{m+1}}$, so the jet factors through
neither branch.
\end{proof}

\begin{Remark}
At the level of genuine arcs, this phenomenon disappears.  If $x(t)y(t)=0$
in $K[[t]]$, then $x(t)=0$ or $y(t)=0$, since $K[[t]]$ is an integral domain.
This explains why the branch--conductor presentation is formulated for
reduced infinite arcs and not for all finite Carter tangent jet schemes.
\end{Remark}

\begin{Theorem}
\label{thm:formal-rep}
Let \((W,\cG)\) be a logarithmic simple adapted co-rank one foliation on a
smooth threefold. Let
$
\Sigmatan:=\Sing\cG\cap |\Dinv|.
$
Then the reduced tangential arc functor centred on \(\Sigmatan\) is
identified with the reduced arc functor of the invariant separatrix divisor
along \(\Sigmatan\). Equivalently,
\[
J_\infty^{\tan}(W,\cG;\Sigmatan)_{\red}
\simeq
J_\infty(\Dinv;\Sigmatan)_{\red}
\]
as reduced pro-functors. No scheme-theoretic equality of non-reduced finite
jet schemes is asserted.
\end{Theorem}

\begin{proof}
The statement is local on \(\Sigmatan\).  In every logarithmic simple adapted
chart, Theorem~\ref{thm:carter-representability} identifies tangential arcs
with arcs on \(\Dinv\) after passage to reduced arc functors.  This is
exactly the asserted identification of the reduced tangential arc functor
centred on \(\Sigmatan\).
\end{proof}
\begin{Lemma}
\label{lem:arcs-pushout}
Let $\sS_W^{\sn}$ be the seminormal branch--conductor pushout.  For every
field extension $K/\bC$, every morphism $\Spec K[[t]]\to \sS_W^{\sn}$
is represented by an arc on at least one normalised branch.  If it has more
than one branch representation, the ambiguity is generated by arcs on the
normalised conductor strata.  More precisely, on the reduced arc functors
relevant to this article there is a coequaliser diagram
\[
\bigsqcup_{\alpha<\beta}J_\infty(C_{\alpha\beta}^\nu)_{\red}
\rightrightarrows
\bigsqcup_\alpha J_\infty(S_\alpha^\nu)_{\red}
\longrightarrow
J_\infty(\sS_W^{\sn})_{\red}.
\]
The two arrows identify branch representatives obtained from the same arc on
a normalised conductor.  No finite-jet analogue is required subsequently.
\end{Lemma}
\begin{proof}
The assertion is local on the simple normal crossing model.  It is enough to
treat the standard affine charts, since the general case is obtained by
adjoining smooth parameters.
First consider the double crossing
$
D=(xy=0)\subset \Spec K[[x,y]].
$
A \(K[[t]]\)-arc on \(D\) is given by two series \(x(t),y(t)\in K[[t]]\)
such that
$
x(t)y(t)=0.
$
Since \(K[[t]]\) is an integral domain, either \(x(t)=0\) or \(y(t)=0\).
Hence every arc on \(D\) factors through at least one of the two branches
$
S_x=(x=0), $ and $ S_y=(y=0).
$
If exactly one of \(x(t)\) and \(y(t)\) is zero, the branch representative is
unique.  If both vanish, then the arc factors through the conductor
$
C=(x=y=0),
$ 
and the two branch representatives are the two images of the same conductor
arc under
\[
C\longrightarrow S_x,
\qquad
C\longrightarrow S_y.
\]
Thus, in the double crossing chart, the reduced arc functor of \(D\) is the
coequaliser of the two conductor maps
\[
J_\infty(C)_{\red}
\rightrightarrows
J_\infty(S_x)_{\red}\sqcup J_\infty(S_y)_{\red}.
\]

Now, consider the triple crossing
$
D=(xyz=0)\subset \Spec K[[x,y,z]].$
A \(K[[t]]\)-arc on \(D\) is given by \(x(t),y(t),z(t)\in K[[t]]\) with
\[
x(t)y(t)z(t)=0.
\]
Again, since \(K[[t]]\) is an integral domain, at least one of the three
series is zero.  Hence the arc factors through at least one branch among
\[
S_x=(x=0),\qquad S_y=(y=0),\qquad S_z=(z=0).
\]
The set of branch representatives is determined by the subset
\[
I(\gamma)=\{i\in\{x,y,z\}: i(t)=0\}.
\]

If \(|I(\gamma)|=1\), the representative is unique.  If
\(|I(\gamma)|=2\), the arc factors through the corresponding pairwise
conductor.  For example, if \(x(t)=y(t)=0\), then \(\gamma\) factors through
\[
C_{xy}=(x=y=0),
\]
and the two representatives on \(S_x\) and \(S_y\) are identified by the two
maps \(C_{xy}\to S_x\) and \(C_{xy}\to S_y\).

If \(|I(\gamma)|=3\), then the arc factors through the triple stratum
$
T=(x=y=z=0).
$
Its images in the three pairwise conductors \(C_{xy},C_{xz},C_{yz}\) have the
same specialisation on \(T\).  The three branch representatives are therefore
identified by the transitive equivalence relation generated by the pairwise
conductor maps.  Thus no additional generator is needed at the triple point:
the pairwise conductor identifications already identify all representatives
of an arc contained in \(T\).
Consequently, in the triple crossing chart, every arc is represented by an
arc on at least one branch, and any two branch representatives of the same
arc differ by a chain of identifications coming from arcs on the pairwise
conductors.  This gives the required coequaliser description.

The same argument applies to an arbitrary SNC divisor
$
D=(x_1\cdots x_r=0)
$
after adjoining smooth parameters.  A \(K[[t]]\)-arc satisfies
\[
x_1(t)\cdots x_r(t)=0,
\]
hence at least one \(x_i(t)\) is zero.  Therefore the arc factors through at
least one branch \(S_i=(x_i=0)\).  If it factors through more than one branch,
the ambiguity is generated by the corresponding pairwise conductor strata
\(S_i\cap S_j\).  Passing to normalisations gives the stated diagram
\[
\bigsqcup_{\alpha<\beta}J_\infty(C_{\alpha\beta}^{\nu})_{\red}
\rightrightarrows
\bigsqcup_\alpha J_\infty(S_\alpha^\nu)_{\red}
\longrightarrow
J_\infty(\sS_W^{\sn})_{\red}.
\]

Finally, the seminormal structure sheaf is constructed as the equaliser of
the branch and conductor rings in \eqref{eq:sn-equaliser-operational}.  Since
\(\Spec K[[t]]\) is reduced and seminormal, morphisms from \(\Spec K[[t]]\)
to the seminormal pushout are exactly compatible branch morphisms modulo the
conductor identifications described above.  This proves the asserted
coequaliser statement for the reduced arc functor.
\end{proof}
\begin{Proposition}
\label{prop:tangential-reduction-diagram}
On every adapted model $W$ there is a coequaliser diagram of reduced pro-functors
\[
\bigsqcup_{\alpha<\beta}J_\infty(C_{\alpha\beta}^\nu)_{\red}
\rightrightarrows
\bigsqcup_{\alpha}J_\infty(S_\alpha^\nu)_{\red}
\longrightarrow J_\infty(\sS_W^{\sn})_{\red}.
\]
Moreover the finite morphism $\sS_W^{\sn}\to \Dinv\subset W$ induces a natural morphism
\[
J_\infty(\sS_W^{\sn})_{\red}
\longrightarrow
J_\infty^{\tan}(W,\cG;\Sigmatan)_{\red}.
\]
The first two arrows are induced by the two maps $C_{\alpha\beta}^\nu\to S_\alpha^\nu$ and $C_{\alpha\beta}^\nu\to S_\beta^\nu$.
\end{Proposition}

\begin{proof}
The structure sheaf of $\sS_W^{\sn}$ is obtained from the branch and conductor rings by the equaliser formula \eqref{eq:sn-equaliser-operational}.  Passing to arcs with integral source reverses this point of view: an arc on the seminormal union is represented by an arc on a branch, and two such representatives are identified precisely when they come from the same arc on a normalised conductor.  Lemma~\ref{lem:arcs-pushout} gives the asserted coequaliser of reduced arc functors.
The morphism $\sS_W^{\sn}\to \Dinv\subset W$ gives a morphism of reduced arc functors
\[
J_\infty(\sS_W^{\sn})_{\red}\longrightarrow J_\infty(\Dinv)_{\red}.
\]
By Theorem~\ref{thm:formal-rep}, the target is identified with the reduced tangential arc sector centred on $\Sigmatan$.  This yields the displayed morphism to $J_\infty^{\tan}(W,\cG;\Sigmatan)_{\red}$.  It is the natural morphism induced by the finite map to $\Dinv$; no closed immersion of arc functors is asserted.
\end{proof}


\begin{Corollary}
\label{cor:branch-decomp}
In a simple adapted chart,
\[
J_\infty^{\tan}(W,\cG;\Sigmatan)_{\red}=\bigcup_{i=1}^r J_\infty(S_i;\Sigmatan)_{\red}.
\]
The branch-intersection locus is
\[
\mathcal I=\bigcup_{i<j}J_\infty(S_i\cap S_j;\Sigmatan)_{\red}.
\]
Every irreducible closed tangential cylinder not contained in $\mathcal I$ has a unique generic branch.
\end{Corollary}

\begin{proof}
Theorem~\ref{thm:formal-rep} identifies the reduced tangential arc functor with the reduced arc functor of the union $\Dinv=\bigcup S_i$.  Since this is a finite union of closed subschemes, every irreducible closed subset is contained in one component.  Outside the branch-intersection locus an arc contained in $S_i$ is not contained in any other $S_j$, and hence the generic branch is unique.
\end{proof}

\section{Local branch reduction and coefficient-one cancellation}
\label{sec:branch-local}

Fix a simple adapted chart and a branch $S_i$.  Put
\[
B_i=\sum_{j\ne i}B_{ij}, \qquad B_{ij}=S_j|_{S_i},
\]
and, for $\Delta_W=\sum a_kT_k$ with $T_k$ transverse to the foliation,
\[
\Theta_i=B_i+\sum_k a_kT_k|_{S_i}.
\]

\begin{Proposition}
\label{prop:branch-id}
The principal branch sector
\[
\cB_i^\circ(W,\cG;\Sigma)=J_\infty(S_i;\Sigma)\setminus \bigcup_{j\ne i}J_\infty(S_i\cap S_j)
\]
is canonically identified with
\[
\{\alpha\in J_\infty(S_i):\alpha(0)\in \Sigma\cap S_i,\ \alpha(\eta)\notin \bigcup_{j\ne i}B_{ij}\}.
\]
On the branch--conductor locus, where $\Sigma\cap S_i=\bigcup_{j\ne i}B_{ij}$, this is the usual condition $\alpha(0)\in B_i$ and $\alpha(\eta)\notin B_i$.
Under this identification, contact with $S_j$ and $T_k$ becomes ordinary contact with $B_{ij}$ and $T_k|_{S_i}$.
\end{Proposition}

\begin{proof}
Let $\alpha:\Spec K[[t]]\to W$ be an arc belonging to the left-hand side. By definition of the $i$-th branch sector, $\alpha$ factors through $S_i$, satisfies $\alpha(0)\in\Sigma\cap S_i$, and its generic point is not contained in any other invariant branch $S_j$, $j\ne i$. Since $S_i\hookrightarrow W$ is a closed immersion, factorisation through $S_i$ is equivalent to a unique arc $\alpha_i:\Spec K[[t]]\to S_i$ such that $\alpha$ is the composite $\Spec K[[t]]\xrightarrow{\alpha_i}S_i\hookrightarrow W$. The condition that the generic point of $\alpha$ does not factor through another $S_j$ is precisely the condition that the generic point of $\alpha_i$ is not contained in $B_{ij}:=S_i\cap S_j$ for any $j\ne i$. Thus the left-hand side identifies with the arcs $\alpha_i$ on $S_i$ satisfying $\alpha_i(0)\in\Sigma\cap S_i$ and $\alpha_i(\eta)\notin\bigcup_{j\ne i}B_{ij}$, where $\eta$ is the generic point of $\Spec K[[t]]$.

Conversely, every arc $\alpha_i:\Spec K[[t]]\to S_i$ with $\alpha_i(0)\in\Sigma\cap S_i$ and $\alpha_i(\eta)\notin\bigcup_{j\ne i}B_{ij}$ gives, by composition with $S_i\hookrightarrow W$, an arc in the $i$-th branch sector. In the branch--conductor charts used for discrepancy computations, the condition $\alpha_i(0)\in\Sigma\cap S_i$ records the invariant traces $\bigcup_{j\ne i}B_{ij}$ and also any additional singular points of the foliation lying on a smooth invariant branch.

It remains to compare contact data.  Let $f_j$ be a local equation of $S_j$ in $W$ and let $g_k$ be a local equation of a transverse boundary component $T_k$. On $S_i$, the restrictions $f_j|_{S_i}$ and $g_k|_{S_i}$ define respectively $B_{ij}$ and $T_k|_{S_i}$. Hence $\ord_\alpha(S_j)=\ord_t(f_j\circ\alpha)=\ord_t((f_j|_{S_i})\circ\alpha_i)=\ord_{\alpha_i}(B_{ij})$ and $\ord_\alpha(T_k)=\ord_t(g_k\circ\alpha)=\ord_t((g_k|_{S_i})\circ\alpha_i)=\ord_{\alpha_i}(T_k|_{S_i})$. Thus the branch-sector arcs and their contact orders are exactly the ordinary arcs and contact orders on $S_i$.
\end{proof}

\begin{Proposition}
\label{prop:monomial-codim}
Let $C_i({\bf m},{\bf n})$ be the cylinder in $\cB_i^\circ$ defined by
$
\ord(B_{ij})=m_j \quad (j\ne i),
 $ and  $
\ord(T_k|_{S_i})=n_k.
$
Then
\[
\codim_{J_\infty(S_i)}C_i({\bf m},{\bf n})=\sum_{j\ne i}m_j+\sum_k n_k
\]
and
\[
\lcodim_{(S_i,\Theta_i)}C_i({\bf m},{\bf n})=
\sum_k(1-a_k)n_k.
\]
\end{Proposition}

\begin{proof}
Let \(S_i\) be the chosen branch and argue at the generic point of the corresponding stratum. As \(S_i\) is smooth in the adapted logarithmic chart, local parameters may be chosen \(u,v\) on \(S_i\) such that the invariant traces and transverse boundary traces which meet the stratum are coordinate curves. Thus, after relabelling, each \(B_{ij}\) or \(T_k|_{S_i}\) is locally given by either \(u=0\) or \(v=0\), and the relevant divisors have simple normal crossings.

Let \(\alpha:\Spec K[[t]]\to S_i\) be an arc and write \(u(t)=u\circ\alpha\), \(v(t)=v\circ\alpha\). The condition \(\ord_\alpha(u)\ge m\) is equivalent to the vanishing of the first \(m\) coefficients of \(u(t)\); hence it imposes \(m\) independent linear conditions on the coefficients of the arc. The same holds for \(v\). Since the divisors are coordinate curves with normal crossings, the imposed conditions are independent. It follows that the ordinary codimension of the contact cylinder is
\[
\sum_{j\ne i}m_j+\sum_k n_k,
\]
where \(m_j=\ord_\alpha(B_{ij})\) and \(n_k=\ord_\alpha(T_k|_{S_i})\), with the evident convention that only the components appearing in the chosen local chart are included.
The branch boundary has the form
\[
\Theta_i=\sum_{j\ne i}B_{ij}+\sum_k a_k\,T_k|_{S_i}.
\]
Therefore the logarithmic codimension of the same cylinder with respect to the pair \((S_i,\Theta_i)\) is
\[
\begin{aligned}
\lcodim_{(S_i,\Theta_i)}(C)
&=\codim_{J_\infty(S_i)}(C)-\ord_C(\Theta_i)\\
&=\left(\sum_{j\ne i}m_j+\sum_k n_k\right)
-\left(\sum_{j\ne i}m_j+\sum_k a_kn_k\right).
\end{aligned}
\]
Thus
\[
\lcodim_{(S_i,\Theta_i)}(C)
=
\sum_k(1-a_k)n_k.
\]
In particular, the invariant contact terms \(m_j\) cancel exactly because every invariant trace \(B_{ij}\) occurs in \(\Theta_i\) with coefficient \(1\).
\end{proof}

\begin{Corollary}
\label{cor:coeff-one}
Contacts with invariant separatrices do not contribute to the weighted tangential codimension.  Only transverse boundary coefficients remain.
\end{Corollary}

\section{The normalised separatrix--conductor system}
\label{sec:system}

\begin{Construction}
Let $\Dinv=\sum S_\alpha$ on an adapted model $W$.  Let
$
\nu_\alpha:S_\alpha^\nu\to S_\alpha
$
be the normalisation.  For two adjacent branches, define the reduced conductor curve
$
C_{\alpha\beta}=S_\alpha\cap S_\beta
$
and let $C_{\alpha\beta}^\nu$ be its normalisation.  The maps
\[
C_{\alpha\beta}^\nu\to S_\alpha^\nu,
\qquad
C_{\alpha\beta}^\nu\to S_\beta^\nu
\]
are induced by the universal property of normalisation.  The \emph{normalised separatrix system} is the finite diagram consisting of all $S_\alpha^\nu$ and all $C_{\alpha\beta}^\nu$ with these incidence maps.
\end{Construction}



\begin{Theorem}
\label{thm:system-exists}
Let $(W,\cG,\Delta_W)$ be an object of $\Adm(X,\cF,\Delta;W_0)$.  The invariant separatrices determine a finite normalised separatrix system and a seminormal separatrix space $\sS^{\sn}$.  The generic arc of every irreducible divisorial tangential cylinder is represented either on a unique normalised branch $S_\alpha^\nu$ or on a unique normalised conductor stratum $C_{\alpha\beta}^\nu$ after removing lower-dimensional specialisations.
\end{Theorem}

\begin{proof}
The preceding construction defines the branch--conductor system. It is shown to represent the adapted toroidal tangential divisorial cylinders.

Let \(C\subset J_\infty^{\tan}(W,\mathcal G;\Sigma_{\tan})_{\red}\) be an irreducible adapted toroidal tangential divisorial cylinder, and let \(\alpha_C:\Spec K[[t]]\to W\) denote its generic arc. By Corollary~\ref{cor:branch-decomp}, the arc \(\alpha_C\) factors through \(\Dinv\). Thus the generic point \(\alpha_C(\eta)\), where \(\eta\) is the generic point of \(\Spec K[[t]]\), has a well-defined generic stratum in the simple normal crossing divisor \(\Dinv\).

If \(\alpha_C(\eta)\) is contained in exactly one component \(S_\alpha\) of \(\Dinv\), then \(\alpha_C\) factors through \(S_\alpha\) and is generically contained in the normal locus of \(S_\alpha\). Since \(\nu_\alpha:S_\alpha^\nu\to S_\alpha\) is the normalisation, the valuative criterion for normalisation gives a unique lift
\[
\begin{tikzcd}
& S_\alpha^\nu \arrow[d,"\nu_\alpha"]\\
\Spec K[[t]] \arrow[r,"\alpha_C"'] \arrow[ur,dashed,"\widetilde\alpha_C"]&
S_\alpha.
\end{tikzcd}
\]
Thus the cylinder \(C\) is represented on the branch \(S_\alpha^\nu\).
If \(\alpha_C(\eta)\) lies on exactly two components, say
\(S_\alpha\cap S_\beta\), and not in a triple intersection, then
\(\alpha_C\) factors through the conductor curve
\(C_{\alpha\beta}=S_\alpha\cap S_\beta\). Since the generic point of the arc lies in the normal locus of \(C_{\alpha\beta}^\nu\), the same valuative criterion gives a unique lift
\[
\begin{tikzcd}
& C_{\alpha\beta}^\nu \arrow[d]\\
\Spec K[[t]] \arrow[r,"\alpha_C"'] \arrow[ur,dashed,"\widetilde\alpha_C"]&
C_{\alpha\beta}.
\end{tikzcd}
\]
Hence the cylinder is represented on \(C_{\alpha\beta}^\nu\).

It remains to consider the case where \(\alpha_C(\eta)\) lies at a triple point
\(p=S_\alpha\cap S_\beta\cap S_\gamma\). Since \(C\) is an adapted toroidal divisorial cylinder, it is determined by a divisorial valuation obtained from a toroidal blow-up of the local coordinate cone. Blow up the triple stratum:
\[
\mu:W'\longrightarrow W.
\]
In adapted local coordinates \(\Dinv=(xyz=0)\), this is the star subdivision of the cone
\[
\sigma=\mathbb R_{\geq 0}\langle e_x,e_y,e_z\rangle
\]
by the ray \(e_x+e_y+e_z\). The lift of the toroidal divisorial datum to \(W'\) has centre either on a unique transformed branch or on a pairwise conductor curve of the refined divisor \(\Dinv'\). Therefore, after this adapted blow-up, the preceding two cases apply to the strict transform of \(C\).

The only arcs remaining entirely at the triple point are constant arcs at \(p\). These form a lower-dimensional specialisation of the cylinders obtained above and do not determine a new divisorial tangential valuation. Hence they do not contribute an additional adapted toroidal tangential divisorial cylinder.

Consequently every adapted toroidal tangential divisorial cylinder is represented either on some normalised branch \(S_\alpha^\nu\) or on some normalised conductor curve \(C_{\alpha\beta}^\nu\), after the allowed adapted toroidal refinement. This is exactly the asserted representation by the branch--conductor system.
\end{proof}


\section{Functoriality under adapted blow-ups}
\label{sec:functoriality}

\begin{Lemma}
\label{lem:branch-transform}
Let $\mu:W'\to W$ be the blow-up of a smooth coordinate stratum of $\Dinv+\Supp\Delta_W$. If the centre is tangent to $\cG$, then the exceptional divisor is invariant for the transformed foliation. Moreover, the invariant divisor on $W'$ is the reduced total transform of $\Dinv$, with the exceptional component included.
\end{Lemma}

\begin{proof}
The statement is local at the centre. Choose an adapted logarithmic coordinate chart $W=\Spec k[[x_1,x_2,x_3]]$ in which $\Dinv$ is a union of coordinate hyperplanes. After relabelling, write $\Dinv=(x_1\cdots x_r=0)$, with $1\le r\le 3$, and assume that the centre is the coordinate stratum $Z=(x_1=\cdots=x_s=0)$, possibly after adjoining transverse boundary coordinates among the remaining variables. In the logarithmic simple adapted chart the foliation is generated by a logarithmic one-form of the shape $\omega=u\,x_1\cdots x_r\eta$, where $u$ is a unit and
\[
\eta=\sum_{i=1}^r\lambda_i\,dx_i/x_i+\sum_{j=r+1}^3 h_j\,dx_j.
\]
The non-dicritical tangent condition for the centre means that the logarithmic residue of $\eta$ in the exceptional direction is nonzero.

It suffices to compute on one standard blow-up chart. On the chart where $x_1$ is exceptional, write $x_i=x_1y_i$ for $2\le i\le s$ and keep the remaining coordinates unchanged. Then $E=(x_1=0)$ and, for $2\le i\le s$, one has $dx_i/x_i=dx_1/x_1+dy_i/y_i$. Therefore the pull-back of the logarithmic part is
\[
\mu^*\eta=\left(\sum_{i=1}^{\min(r,s)}\lambda_i\right)dx_1/x_1+\sum_{2\le i\le \min(r,s)}\lambda_i\,dy_i/y_i+\sum_{i>s,\ i\le r}\lambda_i\,dx_i/x_i+\text{regular terms}.
\]

By the tangent non-dicritical assumption, the coefficient $\sum_{i=1}^{\min(r,s)}\lambda_i$ of $dx_1/x_1$ is nonzero. Hence the transformed logarithmic form has a logarithmic component along $E=(x_1=0)$ with nonzero residue. Thus $E$ is invariant for the transformed foliation.

The strict transforms of the old invariant components are the coordinate hyperplanes $y_i=0$ for $2\le i\le r$ with $i\le s$, together with $x_i=0$ for $i>s$ and $i\le r$. The exceptional component is $x_1=0$. Hence on this chart the invariant divisor is precisely the reduced divisor supported on the total transform of $\Dinv$ together with $E$.
The same computation on the other standard blow-up charts gives the identical conclusion, with the exceptional coordinate replacing $x_1$. Therefore no dicritical exceptional component is produced. Since the assertion is local on the blow-up charts, the invariant divisor on $W'$ is the reduced total transform of $\Dinv$, including the exceptional divisor.
\end{proof}


\begin{Lemma} \label{lem:coordinate-pullback}
Let the logarithmic part of the foliation be
\[
\omega=u_0\,x_1\cdots x_r\left(\sum_{i=1}^r\lambda_i\frac{dx_i}{x_i}+\eta\right),
\]
where $\eta$ is regular in the remaining coordinates.  The following explicit formulas hold.
\begin{enumerate}[label=\textup{(\alph*)},leftmargin=2.4em]
\item If $C=(x=y=0)$ is a conductor curve and $\mu$ is its blow-up, then on the $x$-chart, with $y=xu$,
\begin{equation}\label{eq:curve-pullback}
\mu^*\omega
=u_1x^2u\left((\lambda_x+\lambda_y)\frac{dx}{x}+\lambda_y\frac{du}{u}+\eta_x\right).
\end{equation}
Thus the reduced transformed invariant divisor in this chart is $xu=0$: the exceptional branch is $x=0$ and the strict transform of $S_y$ is $u=0$.
\item If $p=(x=y=z=0)$ is a triple point and $\mu$ is its blow-up, then on the $x$-chart, with $y=xu$ and $z=xv$,
\begin{equation}\label{eq:triple-pullback}
\mu^*\omega
=u_2x^3uv\left((\lambda_x+\lambda_y+\lambda_z)\frac{dx}{x}+\lambda_y\frac{du}{u}+\lambda_z\frac{dv}{v}+\eta_x\right).
\end{equation}
Thus the reduced transformed invariant divisor is $xuv=0$.
\end{enumerate}
\end{Lemma}

\begin{proof}
The two formulas are proved in the indicated affine charts.
For the blow-up of the conductor curve $(x=y=0)$, work on the chart
$y=xu$. Then

\[
\mu^*\left(\frac{dx}{x}\right)=\frac{dx}{x},
\qquad
\mu^*\left(\frac{dy}{y}\right)
=
\frac{d(xu)}{xu}
=
\frac{dx}{x}+\frac{du}{u}.
\]
Therefore
\[
\mu^*\left(
\lambda_x\frac{dx}{x}+\lambda_y\frac{dy}{y}
\right)
=
(\lambda_x+\lambda_y)\frac{dx}{x}
+
\lambda_y\frac{du}{u}.
\]
Moreover,
$
\mu^*(xy)=x(xu)=x^2u.
$
Thus the pulled-back logarithmic generator has the form
\[
x^2u\left(
(\lambda_x+\lambda_y)\frac{dx}{x}
+
\lambda_y\frac{du}{u}
+
\text{terms in the remaining variables}
\right).
\]
After saturating the pulled-back foliation by the common divisorial factor,
the logarithmic residues on this chart are
$
\lambda_E=\lambda_x+\lambda_y,
 $ and $
\lambda_{S_y'}=\lambda_y.
$
The exceptional divisor is $E=(x=0)$ and the strict transform of $(y=0)$ is
$(u=0)$. Thus \eqref{eq:curve-pullback} follows on this chart. The other
chart $x=yv$ is identical, with $x$ and $y$ interchanged.
For the blow-up of the triple point $(x=y=z=0)$, work on the chart
$
y=xu, z=xv.
$
Then
\[
\mu^*\left(\frac{dx}{x}\right)=\frac{dx}{x},
\qquad
\mu^*\left(\frac{dy}{y}\right)=\frac{dx}{x}+\frac{du}{u},
\qquad
\mu^*\left(\frac{dz}{z}\right)=\frac{dx}{x}+\frac{dv}{v}.
\]
Hence
\[
\mu^*\left(
\lambda_x\frac{dx}{x}
+\lambda_y\frac{dy}{y}
+\lambda_z\frac{dz}{z}
\right)
=
(\lambda_x+\lambda_y+\lambda_z)\frac{dx}{x}
+
\lambda_y\frac{du}{u}
+
\lambda_z\frac{dv}{v}.
\]
Also,
$
\mu^*(xyz)=x(xu)(xv)=x^3uv.
$
Therefore the pulled-back logarithmic generator is
\[
x^3uv\left(
(\lambda_x+\lambda_y+\lambda_z)\frac{dx}{x}
+
\lambda_y\frac{du}{u}
+
\lambda_z\frac{dv}{v}
+
\text{regular terms}
\right).
\]
After saturation by the common divisorial factor, the residues along the new
coordinate invariant components are
\[
\lambda_E=\lambda_x+\lambda_y+\lambda_z,
\qquad
\lambda_{S_y'}=\lambda_y,
\qquad
\lambda_{S_z'}=\lambda_z.
\]
Thus \eqref{eq:triple-pullback} follows on this chart. The remaining affine
charts are obtained by choosing $y$ or $z$ as the exceptional coordinate and
give the same residue rule with the variables permuted.

In both cases the positive non-resonance condition of Definition
\ref{def:simple-chart} implies that the sums of residues appearing in the
exceptional direction are non-zero for the adapted blow-ups considered in this article.
Consequently the exceptional component is invariant exactly in these
tangent non-dicritical charts, and the invariant components after blow-up
are precisely the displayed coordinate components.
\end{proof}

\begin{center}
\begin{tikzcd}[column sep=large,row sep=large]
S_{\alpha}'^{\nu} \arrow[d] \arrow[r] & S_{\alpha}^{\nu} \arrow[d] \\
\sS^{\sn}_{W'} \arrow[r] & \sS^{\sn}_{W}
\end{tikzcd}
\qquad
\begin{tikzcd}[column sep=large,row sep=large]
E \arrow[d] \arrow[r] & C_{\alpha\beta}^{\nu} \arrow[d] \\
\sS^{\sn}_{W'} \arrow[r] & \sS^{\sn}_{W}
\end{tikzcd}
\end{center}

\begin{Proposition}
\label{prop:branch-functoriality}
Let $\mu:W'\to W$ be an adapted blow-up.  The normalised branch--conductor diagram on $W'$ maps to the diagram on $W$ in the following sense.
\begin{enumerate}[label=\textup{(\roman*)},leftmargin=2.4em]
\item The strict transform of an old branch maps birationally to that branch.
\item If the centre is a conductor curve $C_{\alpha\beta}$, the exceptional invariant branch is a ruled surface over $C_{\alpha\beta}$; it is a new branch of the refined system, and its structure morphism factors through the old normalised conductor.  It is not birational to the conductor curve.
\item If the centre is a triple point, the blow-up is the toric star subdivision of the local coordinate cone.  The exceptional branch and the strict transforms carry an incidence diagram whose pairwise conductors map to the old pairwise conductors or to the old triple point.
\end{enumerate}
These maps are compatible with the incidence relations and induce a morphism of seminormal separatrix spaces
\[
\sS^{\sn}_{W'}\to \sS^{\sn}_{W}.
\]
\end{Proposition}

\begin{proof}
The strict transform case follows from the functoriality of normalisation. Indeed, if \(S_\gamma'\) is the strict transform of \(S_\alpha\), then \(\mu(S_\gamma')\subset S_\alpha\). Hence the universal property of normalisation gives a unique morphism \(S_\gamma^{\prime\nu}\to S_\alpha^\nu\) fitting into the commutative diagram
\[
\begin{tikzcd}
S_\gamma^{\prime\nu} \arrow[r] \arrow[d] &
S_\gamma' \arrow[d,"\mu|_{S_\gamma'}"]\\
S_\alpha^\nu \arrow[r] &
S_\alpha.
\end{tikzcd}
\]

Suppose now that the centre is the conductor curve \(C_{\alpha\beta}=S_\alpha\cap S_\beta\). Choose local coordinates with \(S_\alpha=(x=0)\) and \(S_\beta=(y=0)\). Thus \(\mu\) is locally the blow-up of the ideal \((x,y)\). The exceptional divisor is \(E=\mathbb P(N_{C_{\alpha\beta}/W})\), and the natural projection \(E\to C_{\alpha\beta}\) induces, after normalisation, a morphism \(E^\nu\to C_{\alpha\beta}^\nu\). Therefore the exceptional component is a new branch of the refined system, and its structural morphism to the old system factors as
\[
E^\nu\longrightarrow C_{\alpha\beta}^\nu\longrightarrow \mathscr S_W^{\sn}.
\]
This is the required functorial behaviour: \(E\) maps to the old conductor stratum; it is not identified birationally with the curve \(C_{\alpha\beta}\).
The new conductor curves are \(E\cap S_\alpha'\) and \(E\cap S_\beta'\). Their maps to the old system both factor through \(C_{\alpha\beta}^\nu\). Moreover, the restriction maps from \(E^\nu\), \(S_\alpha^{\prime\nu}\), and \(S_\beta^{\prime\nu}\) to these new conductor curves are compatible with the two old restrictions to \(C_{\alpha\beta}^\nu\). Equivalently, the following diagrams commute:
\[
\begin{tikzcd}
(E\cap S_\alpha')^\nu \arrow[r] \arrow[d] &
E^\nu \arrow[d]\\
C_{\alpha\beta}^\nu \arrow[r] &
\mathscr S_W^{\sn}
\end{tikzcd}
\qquad
\begin{tikzcd}
(E\cap S_\alpha')^\nu \arrow[r] \arrow[d] &
S_\alpha^{\prime\nu} \arrow[d]\\
C_{\alpha\beta}^\nu \arrow[r] &
S_\alpha^\nu
\end{tikzcd}
\]
and similarly with \(\alpha\) replaced by \(\beta\). Hence the blow-up along a conductor curve defines a morphism of branch--conductor diagrams.

It remains only to treat a triple point. Choose adapted coordinates in which the invariant divisor is \(xyz=0\). Blowing up the origin gives the toric star subdivision of the cone \(\mathbb R_{\geq 0}\langle e_x,e_y,e_z\rangle\) by the ray \(e_x+e_y+e_z\). On the affine chart \(y=xu\), \(z=xv\), the transformed invariant divisor is the coordinate arrangement \(xuv=0\); the other charts are obtained by permuting \(x,y,z\). Thus the refined branch--conductor incidence diagram is the subdivision of the old coordinate incidence diagram.

Each new branch and each new conductor has a canonical image stratum in the old diagram. The exceptional branch maps to the old triple point. A new conductor curve contained in the intersection of two transformed coordinate components maps to the corresponding old conductor curve when its edge lies over an old two-dimensional face, and otherwise maps to the old triple point. These assignments are induced by \(\mu:W'\to W\), and they respect all conductor inclusions.
Consequently, in all cases, one obtains a commutative diagram of branch--conductor systems
\[
\begin{tikzcd}
\displaystyle\bigsqcup_{\gamma<\delta} C_{\gamma\delta}^{\prime\nu}
\arrow[r,shift left=.6ex] \arrow[r,shift right=.6ex] \arrow[d] &
\displaystyle\bigsqcup_{\gamma} S_\gamma^{\prime\nu}
\arrow[d]\\
\displaystyle\bigsqcup_{\alpha<\beta} C_{\alpha\beta}^{\nu}
\arrow[r,shift left=.6ex] \arrow[r,shift right=.6ex] &
\displaystyle\bigsqcup_{\alpha} S_\alpha^\nu.
\end{tikzcd}
\]
Passing to structure sheaves gives a morphism between the corresponding equaliser diagrams. Therefore a compatible tuple of functions on the old normalised branches pulls back to a compatible tuple of functions on the refined normalised branches. By the equaliser description of the seminormal pushout, this induces a unique morphism
\[
\mathscr S_{W'}^{\sn}\longrightarrow \mathscr S_W^{\sn}.
\]
\end{proof}

\begin{Theorem}
\label{thm:common-refinement}
Any two objects in the toroidal subcategory generated from a fixed logarithmic non-dicritical adapted model of $(X,\cF,\Delta)$ admit a common adapted toroidal refinement.  On this refinement the induced morphisms of normalised separatrix systems are compatible with branch and conductor strata.
\end{Theorem}

\begin{proof}
In the fixed toroidal category, the models under consideration are described locally by regular fans whose cones encode the coordinate strata of \(\Dinv+\Supp\Delta_W\). Let \(\Sigma_1\) and \(\Sigma_2\) be the fans associated with two such adapted toroidal models over \(W_0\). Their common refinement is the fan
\[
\Sigma_1\wedge \Sigma_2
=
\{\sigma_1\cap\sigma_2\mid \sigma_1\in\Sigma_1,\ \sigma_2\in\Sigma_2\}.
\]
After subdividing further if necessary, choose a regular fan \(\Sigma_3\) refining \(\Sigma_1\wedge\Sigma_2\).

Every regular subdivision of a smooth toroidal fan can be obtained as a sequence of star subdivisions along smooth coordinate cones. Geometrically, each such star subdivision is the blow-up of the corresponding smooth coordinate stratum of \(\Dinv+\Supp\Delta_W\). Thus the toroidal morphisms associated with the refinements \(\Sigma_3\to\Sigma_1\) and \(\Sigma_3\to\Sigma_2\) are realised by sequences of blow-ups
\[
W_3\longrightarrow W_1,
\qquad
W_3\longrightarrow W_2
\]
along smooth coordinate strata.
By Lemma~\ref{lem:coordinate-pullback}, in every affine chart of such a blow-up the logarithmic residues transform by replacing the exceptional residue with the sum of the old residues along the coordinates defining the centre. The positive non-resonance condition in the logarithmic simple adapted setting implies that this exceptional residue is non-zero for the adapted centres used in this article. Therefore each elementary blow-up remains logarithmic, non-dicritical, and adapted, and the transformed invariant divisor is the reduced total transform of the old invariant divisor.

For each elementary blow-up, Proposition~\ref{prop:branch-functoriality} gives a morphism of normalised branch--conductor diagrams. Composing these morphisms along the two sequences of blow-ups gives compatible morphisms
\[
\mathfrak D_{W_3}\longrightarrow \mathfrak D_{W_1},
\qquad
\mathfrak D_{W_3}\longrightarrow \mathfrak D_{W_2},
\]
and hence, by the equaliser description of the seminormal pushout, morphisms
\[
\mathscr S_{W_3}^{\sn}\longrightarrow \mathscr S_{W_1}^{\sn},
\qquad
\mathscr S_{W_3}^{\sn}\longrightarrow \mathscr S_{W_2}^{\sn}.
\]
Thus \(W_3\) is a common adapted toroidal refinement of \(W_1\) and \(W_2\), and the branch--conductor systems are functorial with respect to this refinement.
This proves only the toroidal common-refinement statement inside the fixed category generated from \(W_0\).
\end{proof}


\section{Normalised foliated adjunction and conductor compatibility}
\label{sec:adjunction}

\begin{Definition} 
For an invariant branch $S_\alpha$ define $\Theta_\alpha$ on $S_\alpha^\nu$ by
\[
K_{S_\alpha^\nu}+\Theta_\alpha
\sim_\bQ
\nu_\alpha^*((K_\cG+\Delta_W)|_{S_\alpha}).
\]
The class $(K_\cG+\Delta_W)|_{S_\alpha}$ denotes the invariant-divisor adjunction class supplied by foliated adjunction, not the naive restriction of an arbitrary Cartier divisor.
Equivalently, in an SNC chart, $\Theta_\alpha$ is the sum of the coefficient-one traces of the other invariant branches, the restrictions of the transverse boundary components with their original coefficients, and the ordinary normalisation/conductor different.
\end{Definition}

\begin{Theorem}
\label{thm:normalised-adjunction}
Let $S$ be an invariant separatrix branch of a logarithmic simple adapted co-rank one foliation on a smooth threefold.  Then there is a canonically determined boundary $\Theta_S$ on $S^\nu$ such that
\[
\nu^*((K_\cG+\Delta_W)|_S)
\sim_\bQ
K_{S^\nu}+\Theta_S.
\]
In this theorem $(K_\cG+\Delta_W)|_S$ is the invariant-divisor adjunction class of foliated adjunction.
The coefficient of every other invariant trace in $\Theta_S$ is one, the transverse components of $\Delta_W$ keep their coefficients, and the normalisation conductor contributes through the ordinary different.  This construction is compatible with adapted blow-ups.
\end{Theorem}

\begin{proof}
Foliated adjunction is used for invariant divisors of co-rank one foliated pairs in the setting of Cascini--Spicer and Spicer--Svaldi \cite{CSAdj,SS}.  Thus, if $S\subset W$ is invariant and $\nu:S^\nu\to S$ is the normalisation, there is a different on $S^\nu$ such that
\[
\nu^*((K_\cG+\Delta_W)|_S)\sim_\bQ K_{S^\nu}+\Diff_S(\Delta_W).
\]
The class $(K_\cG+\Delta_W)|_S$ denotes the adjunction class supplied by the invariant-divisor case of foliated adjunction, not the naive restriction of a Cartier representative.

In a logarithmic simple adapted chart, let $S=(x_1=0)$ and write $\Dinv=(x_1\cdots x_r=0)$.  If $S_j=(x_j=0)$ is another invariant branch, its trace $B_{1j}=S\cap S_j$ appears in the adjunction boundary with coefficient one.  If $T_k$ is a transverse boundary component with coefficient $a_k$ in $\Delta_W$, then its trace on $S$ appears with coefficient $a_k$.  If $S$ is not normal, the pull-back to $S^\nu$ also contains the ordinary conductor different.  This gives the boundary
\[
\Theta_S=
\sum_{j\ne 1}B_{1j}^\nu+
\sum_k a_k\,\nu^*(T_k|_S)+
\operatorname{Cond}_{S^\nu/S},
\]
and hence
 $
\nu^*((K_\cG+\Delta_W)|_S)\sim_\bQ K_{S^\nu}+\Theta_S.
$
Let $\mu:W'\to W$ be an adapted blow-up and let $S'$ be a transformed branch mapping to $S$.  Write $\sigma:S'^\nu\to S^\nu$ for the induced morphism.  Applying invariant-divisor foliated adjunction on $S$ and on $S'$ gives two expressions for the restriction of the same ambient birational class.  The boundary on $S'^\nu$ is the crepant transform of the boundary on $S^\nu$ together with the boundary terms attached to any new branch or conductor stratum of the refined separatrix system.

For a point centre on a smooth branch surface, let $b$ be the sum of the coefficients of the components of $\Theta_S$ passing through the centre, and let $F_S$ be the exceptional curve of $\sigma$.  Then
\[
K_{S'}=\sigma^*K_S+F_S,
\qquad
\sigma^*\Theta_S=\sigma^{-1}_*\Theta_S+bF_S.
\]
Thus the crepant transform of the boundary is
$
\Theta_{S'}^{\mathrm cr}=\sigma^{-1}_*\Theta_S+(b-1)F_S,
$
and
\[
K_{S'}+\Theta_{S'}^{\mathrm cr}=\sigma^*(K_S+\Theta_S).
\]
This is the surface calculation used in the comparison.  It replaces the erroneous practice of adding a further exceptional term to the crepant adjunction identity.

If the ambient centre is a conductor curve, then its intersection with an old branch is a Cartier divisor on that branch, and the blow-up of the smooth branch along this divisor is an isomorphism.  The ambient exceptional divisor is a new invariant branch of the refined system, and its intersections with the strict transforms of the old branches are new conductor curves.  It is therefore recorded in the refined branch--conductor diagram, not as an exceptional curve on the old branch surface.

These local computations are compatible with the conductor restrictions.  Hence the adjunction boundaries are crepant compatible under adapted blow-ups in the fixed toroidal category.
\end{proof}

\begin{Proposition}
\label{prop:no-double-count}
Let $C_{\alpha\beta}^\nu$ be the normalised conductor curve associated with two adjacent invariant branches.  The boundaries induced on $C_{\alpha\beta}^\nu$ from $(S_\alpha^\nu,\Theta_\alpha)$ and $(S_\beta^\nu,\Theta_\beta)$ agree.  At a triple point the three pairwise restrictions have the same further restriction to the normalised point stratum.
\end{Proposition}

\begin{proof}
At the generic point of \(C_{\alpha\beta}\), choose adapted coordinates
$
S_\alpha=(x=0), S_\beta=(y=0) $ and $ C_{\alpha\beta}=(x=y=0).
$
Let \(T=(z=0)\) be a transverse boundary component with coefficient \(a\). On the branch \(S_\alpha\), the conductor \(C_{\alpha\beta}\) is defined by \(y|_{S_\alpha}=0\), while on the branch \(S_\beta\) it is defined by \(x|_{S_\beta}=0\). By invariant-divisor foliated adjunction, the other invariant branch appears in the boundary with coefficient \(1\), and the transverse component \(T\) appears with its original coefficient \(a\). Hence, near the generic point of \(C_{\alpha\beta}\),
\[
\Theta_\alpha=(y=0)|_{S_\alpha}+a(z=0)|_{S_\alpha}+\Theta_\alpha',
\qquad
\Theta_\beta=(x=0)|_{S_\beta}+a(z=0)|_{S_\beta}+\Theta_\beta',
\]
where no component of \(\Theta_\alpha'\) or \(\Theta_\beta'\) contains the generic point of \(C_{\alpha\beta}\).
Let \(\rho_\alpha:C_{\alpha\beta}^{\nu}\to S_\alpha^\nu\) and \(\rho_\beta:C_{\alpha\beta}^{\nu}\to S_\beta^\nu\) be the two conductor maps. Removing the conductor component itself from the two branch boundaries and restricting to \(C_{\alpha\beta}^{\nu}\), one obtains
\[
\rho_\alpha^*(\Theta_\alpha-C_{\alpha\beta})
=
a(z=0)|_{C_{\alpha\beta}^{\nu}}
=
\rho_\beta^*(\Theta_\beta-C_{\alpha\beta})
\]
away from triple points. Equivalently, the two branch adjunction formulas induce the same adjunction class on the normalised conductor:
\[
\rho_\alpha^*(K_{S_\alpha^\nu}+\Theta_\alpha-C_{\alpha\beta})
\sim_{\mathbb Q}
K_{C_{\alpha\beta}^{\nu}}+\Theta_{\alpha\beta}
\sim_{\mathbb Q}
\rho_\beta^*(K_{S_\beta^\nu}+\Theta_\beta-C_{\alpha\beta}).
\]
Thus the conductor component is not counted twice; it is the gluing stratum for the two branch boundaries.

At a triple point, choose coordinates with \(\Dinv=(xyz=0)\). Consider, for example, the conductor \(C_{xy}=(x=y=0)\). On the branch \(S_x=(x=0)\), the boundary contains the traces \(C_{xy}\) and \(C_{xz}\) with coefficient \(1\). On the branch \(S_y=(y=0)\), it contains \(C_{xy}\) and \(C_{yz}\) with coefficient \(1\). After subtracting the common conductor \(C_{xy}\), the remaining invariant traces \(C_{xz}\) and \(C_{yz}\) restrict to the same reduced closed point \(x=y=z=0\) on \(C_{xy}^{\nu}\). Therefore
\[
(C_{xz})|_{C_{xy}^{\nu}}
=
(x=y=z=0)
=
(C_{yz})|_{C_{xy}^{\nu}}
\]
as reduced zero-dimensional divisors. The same computation holds for the other two pairwise conductors. Hence the restrictions of the branch boundaries agree on the normalised point stratum, and the triple-point contribution is a single gluing datum in the branch--conductor system.
\end{proof}


\section{Tangential divisorial valuations}
\label{sec:valuations}

A minor terminological point is important.  A divisor over a branch surface does not, by itself, define a divisorial valuation of the threefold function field $K(X)$.  What it defines canonically is a divisorial datum in the tangential arc sector.  When this datum is realised by an exceptional divisor on an adapted threefold model, it also gives an ordinary divisorial valuation over $X$.  The subsequent definitions keep these two levels separate.

\begin{Definition}
An \emph{adapted toroidal tangential divisorial datum} is either
\begin{enumerate}[label=\textup{(\roman*)},leftmargin=2.4em]
\item a prime divisor $F$ over a normalised branch surface $S_\alpha^\nu$, or
\item a closed point $p\in C_{\alpha\beta}^\nu$ of a normalised conductor curve, equivalently the order valuation $\ord_p$,
\end{enumerate}
considered up to common adapted refinement and the conductor identifications of Proposition~\ref{prop:no-double-count}.  Such a datum determines a maximal divisorial cylinder in the tangential arc space.  If it is represented by an exceptional divisor on an adapted threefold model, it is also denoted that threefold divisor by $E$.
\end{Definition}

\begin{Theorem}
\label{thm:classification}
Let $E$ be an adapted toroidal tangential divisorial datum whose centre is generically contained in the invariant separatrix locus on some adapted model.  After passing to an adapted refinement, the tangential cylinder determined by $E$ is represented by a divisor over a normalised branch surface or by a conductor datum on a normalised conductor curve.  Conversely, every branch or conductor datum defines a divisorial tangential cylinder, and those which occur as exceptional divisors on adapted threefold models give divisorial valuations over $X$.
\end{Theorem}

 \begin{proof}
The statement is local in a toroidal chart for
\(\Dinv+\Supp\Delta_W\).  Let \(\Sigma_W\) be the local fan and let
\(\sigma\in\Sigma_W\) be the cone corresponding to the coordinate stratum
which contains the centre of the datum.  By definition, an adapted toroidal
tangential divisorial datum is monomial after an adapted toroidal refinement.
Equivalently, it is represented by a primitive vector
\[
v\in N\cap\sigma .
\]

Choose a regular star subdivision of \(\sigma\) which contains the ray
\(\rho_v:=\mathbb R_{\geq0}v\).  This subdivision is realised by a finite
sequence of blow-ups of smooth coordinate strata.  By
Lemma~\ref{lem:coordinate-pullback}, and by positive non-resonance at each
step, these blow-ups are adapted and remain in the logarithmic
non-dicritical setting.  On the resulting model, the ray \(\rho_v\)
corresponds to a prime divisor \(E_v\).  Thus the given toroidal datum is
realised by the divisorial centre of \(E_v\).

We now read this centre on the reduced tangential sector.  By
Theorem~\ref{thm:formal-rep},
\[
J_\infty^{\tan}(W,\cG;\Sigmatan)_{\red}
=
J_\infty(\Dinv;\Sigmatan)_{\red}.
\]
By Lemma~\ref{lem:arcs-pushout}, the right-hand side is presented by the
coequaliser
\[
\bigsqcup_{\alpha<\beta}J_\infty(C_{\alpha\beta}^{\nu})_{\red}
\rightrightarrows
\bigsqcup_{\alpha}J_\infty(S_\alpha^\nu)_{\red}
\longrightarrow
J_\infty(\sS_W^{\sn})_{\red}.
\]
Hence every irreducible tangential divisorial cylinder is represented either
on a normalised branch or on a normalised conductor stratum, with the usual
conductor identifications.
More explicitly, let \(Z_v\) be the generic centre of \(E_v\) on the
separatrix divisor of the refined model.  There are only the following
possibilities:
\[
Z_v\subset S_\alpha^\circ,\qquad
Z_v\subset C_{\alpha\beta}^\circ,\qquad
Z_v\subset S_\alpha\cap S_\beta\cap S_\gamma .
\]
In the first case, the corresponding arcs have generic branch \(S_\alpha\);
by the branch-sector identification they are represented by a maximal
divisorial cylinder in \(J_\infty(S_\alpha^\nu)\).  In the second case, the
arcs are generically supported on the conductor, and the datum is represented
by the corresponding maximal divisorial cylinder in
\(J_\infty(C_{\alpha\beta}^\nu)\).

It remains only to remove the third possibility.  If locally
\(\Dinv=(xyz=0)\), then the cone is
\[
\sigma=\mathbb R_{\geq0}\langle e_x,e_y,e_z\rangle .
\]
A regular star subdivision extracting \(\rho_v\) refines \(\sigma\) into
regular cones whose faces correspond to transformed branches and pairwise
conductors.  After extracting \(\rho_v\), the generic centre of the divisor
\(E_v\) is a codimension-one stratum in the toroidal boundary of the refined
model.  Its intersection with the reduced separatrix system is therefore read
on either a single transformed branch or on a pairwise conductor.  Thus the
triple-stratum case is reduced to one of the two cases above.  The process is
finite because a fixed primitive ray is extracted by a finite regular
subdivision of the cone containing it.

Conversely, let \(F\) be a divisor over a normalised branch
\(S_\alpha^\nu\).  The ordinary maximal divisorial set \(N_q(F)\subset
J_\infty(S_\alpha^\nu)\) maps, through the branch-sector identification and
the coequaliser presentation above, to an adapted tangential divisorial
cylinder.  The same construction applies to a divisor over a normalised
conductor \(C_{\alpha\beta}^\nu\), giving a conductor-type tangential
divisorial cylinder.

If the branch or conductor datum is realised by a prime divisor on an adapted
toroidal threefold model over \(W\), it defines a divisorial valuation of
\(K(W)=K(X)\).  Otherwise it remains a tangential divisorial datum on the
normalised branch--conductor system.  This is exactly the distinction built
into the adapted toroidal tangential sector.
\end{proof}

\begin{Definition}
Let an adapted toroidal tangential divisorial datum be represented by $F$ over a stratum $(V,B_V)$ of the normalised separatrix system.  Define
\[
a_{\tan}(F;X,\cF,\Delta)=a(F;V,B_V).
\]
For conductor type data, take $V=C_{\alpha\beta}^\nu$ and $B_V=\Theta_{\alpha\beta}$.  When the datum is realised by a threefold divisor $E$, write $a_{\tan}(E;X,\cF,\Delta)$.
\end{Definition}

\begin{Proposition}
\label{prop:atan-invariance}
The number $a_{\tan}$ is independent of the chosen representative inside the fixed adapted toroidal category and of the branch presentation.
\end{Proposition}

\begin{proof}
Let two branch--conductor presentations of the same adapted toroidal tangential divisorial datum be given. By Theorem~\ref{thm:common-refinement}, there exists a common adapted toroidal refinement
\[
W_3\longrightarrow W_1,
\qquad
W_3\longrightarrow W_2.
\]
By Proposition~\ref{prop:branch-functoriality}, the normalised branch--conductor diagram of \(W_3\) maps to the normalised branch--conductor diagrams of both \(W_1\) and \(W_2\). Hence the stratum representing the datum on \(W_3\) dominates the two strata representing the datum on \(W_1\) and \(W_2\).

Assume first that the datum is represented on a branch. Let \(S_i^{(a),\nu}\) be the normalised branch on \(W_a\), for \(a=1,2,3\), and let
$
\sigma_a:S_i^{(3),\nu}\longrightarrow S_i^{(a),\nu}
$
be the induced morphisms on the common refinement. By Theorem~\ref{thm:normalised-adjunction}, the branch adjunction boundaries are crepant compatible:
\[
K_{S_i^{(3),\nu}}+\Theta_i^{(3)}
\sim_{\mathbb Q}
\sigma_a^*(K_{S_i^{(a),\nu}}+\Theta_i^{(a)})
\]
over the generic centre of the datum. Therefore, for the divisorial valuation \(E\) corresponding to the datum,
$
a_E(S_i^{(3),\nu},\Theta_i^{(3)})
=
a_E(S_i^{(a),\nu},\Theta_i^{(a)})
$
for \(a=1,2\). Hence the value computed from the two branch presentations agrees.

Now, suppose that the datum is represented in a conductor sector. Let \(C_{\alpha\beta}^{(a),\nu}\) be the normalised conductor stratum on \(W_a\), and let
$
\tau_a:C_{\alpha\beta}^{(3),\nu}\longrightarrow C_{\alpha\beta}^{(a),\nu}
$
be the induced maps from the common refinement. By Proposition~\ref{prop:no-double-count}, the restrictions of the two adjacent branch boundaries to the conductor agree after removing the conductor component itself. Thus the two adjacent branch computations induce one and the same conductor boundary, say \(\Theta_{\alpha\beta}^{(a)}\), on \(C_{\alpha\beta}^{(a),\nu}\). The crepant compatibility of Theorem~\ref{thm:normalised-adjunction} restricts to
\[
K_{C_{\alpha\beta}^{(3),\nu}}+\Theta_{\alpha\beta}^{(3)}
\sim_{\mathbb Q}
\tau_a^*(K_{C_{\alpha\beta}^{(a),\nu}}+\Theta_{\alpha\beta}^{(a)}).
\]
Ordinary log discrepancy is invariant under crepant birational pull-back. Hence
$
a_E(C_{\alpha\beta}^{(3),\nu},\Theta_{\alpha\beta}^{(3)})
=
a_E(C_{\alpha\beta}^{(a),\nu},\Theta_{\alpha\beta}^{(a)})
$
for \(a=1,2\). Therefore the conductor presentation also gives the same value on both models.
The branch and conductor cases exhaust the adapted toroidal tangential divisorial data by the classification theorem. Consequently the value assigned to the datum is independent of the chosen adapted model, the chosen branch--conductor presentation, and the chosen common refinement.
\end{proof}

\section{Local discrepancy calculations}
\label{sec:local-discrepancy}

Let $Z$ be a smooth tangent stratum in a simple chart.  Write
\[
Z=\bigcap_{\lambda\in I}S_\lambda\cap\bigcap_{k\in K}T_k,
\qquad m=|I|,
\qquad \beta_Z=\sum_{k\in K}a_k.
\]
Choose a branch $S_i$ with $i\in I$.

\begin{Proposition}
\label{prop:elementary-blowups}
Let $\mu:W'\to W$ be the blow-up of $Z$ and let $E$ be the exceptional divisor.
\begin{enumerate}[label=\textup{(\roman*)},leftmargin=2.4em]
\item If $\codim_W Z=2$, then
\[
\AF(E;W,\cG,\Delta_W)=1-(m-1)-\beta_Z.
\]
Thus an invariant curve $S_i\cap S_j$ gives discrepancy $0$, while a curve $S_i\cap T_k$ gives $1-a_k$.
\item If $\codim_W Z=3$, then
\[
\AF(E;W,\cG,\Delta_W)=2-(m-1)-\beta_Z.
\]
The three possible values are $2-a_{k_1}-a_{k_2}$, $1-a_k$, and $0$.
\end{enumerate}
These same numbers are the ordinary discrepancies of the induced branch blow-up on $(S_i,\Theta_i)$.
\end{Proposition}

\begin{proof}
Let \(Z\subset W\) be the centre of the adapted toroidal blow-up and assume that the tangential datum is represented on the branch \(S_i\). Let \(m\) be the number of invariant components of \(\Dinv\) containing \(Z\). Then, on the surface \(S_i\), the centre is contained in exactly \(m-1\) invariant traces \(B_{ij}=S_i\cap S_j\), each appearing in the branch boundary with coefficient \(1\). Let
\[
\beta_Z:=\sum_{T_k\supset Z} a_k
\]
be the total coefficient of the transverse boundary components passing through \(Z\). Thus the total boundary coefficient through the centre on \(S_i\) is
$
(m-1)+\beta_Z.
$

If \(Z\) is a curve in \(W\), then \(Z\subset S_i\) is a divisor on the
smooth surface \(S_i\).  The ambient blow-up has an exceptional divisor on the
threefold, but the branch adjunction system reads its tangential contribution
as the divisorial valuation of the boundary curve \(Z\subset S_i\), not as an
exceptional divisor over \(S_i\).  The ordinary log discrepancy of this
divisor with respect to the surface pair \((S_i,\Theta_i)\) is
\[
a_Z(S_i,\Theta_i)
=
1-\operatorname{coeff}_Z(\Theta_i)
=
1-\bigl((m-1)+\beta_Z\bigr).
\]

If \(Z\) is a point in \(W\), then \(Z\) is a point on the surface \(S_i\). Let
$
\sigma:S_i'\longrightarrow S_i
$
be the blow-up of this point, and let \(F\) be the exceptional curve. As \(S_i\) is smooth,
$
K_{S_i'}=\sigma^*K_{S_i}+F.
$
Moreover, the pull-back of the boundary components through \(Z\) contributes
\[
\sigma^*\Theta_i
=
\Theta_i' + \bigl((m-1)+\beta_Z\bigr)F
\]
near the generic point of \(F\). Hence
\[
K_{S_i'}+\Theta_i'
=
\sigma^*(K_{S_i}+\Theta_i)
+
\left(1-\bigl((m-1)+\beta_Z\bigr)\right)F.
\]
Equivalently, the log discrepancy of \(F\) over the surface pair is
\[
a_F(S_i,\Theta_i)
=
2-\bigl((m-1)+\beta_Z\bigr).
\]

On the foliated threefold side, the blow-up is tangent. In the foliated discrepancy convention, invariant components contribute by the coefficient-one rule, equivalently through \(\epsilon(E)=0\) for tangent exceptional divisors, while transverse boundary components subtract their coefficients. Thus the invariant part contributes exactly \(m-1\), and the transverse part contributes \(\beta_Z\). Therefore the foliated tangential discrepancy attached to the adapted blow-up is
$
1-\bigl((m-1)+\beta_Z\bigr)
$
when the centre is a curve and
$
2-\bigl((m-1)+\beta_Z\bigr)
$
when the centre is a point on the branch. These are precisely the ordinary log discrepancies computed on the branch pair \((S_i,\Theta_i)\).
\end{proof}

\section{Comparison with foliated discrepancies}
\label{sec:foliated-comparison}

The tangential discrepancy $a_{\tan}$ is defined on the normalised
branch--conductor adjunction system.  We now compare it with the foliated MMP
discrepancy on the part where such a comparison is intrinsic: adapted
toroidal invariant divisors whose data are read on normalised invariant
branches.  Conductor data remain essential for $\tmldtor$, but they are
adjunction strata rather than prime divisors over the foliated threefold.

\begin{Notation}
We write $\AF(E;X,\cF,\Delta)$ for the foliated log discrepancy of an
invariant divisor $E$, in the convention compatible with the foliated MMP.
\end{Notation}

\begin{Lemma} 
\label{lem:foliated-discrepancy-additivity}
Let
\[
W_N\xrightarrow{\mu_N}W_{N-1}\longrightarrow\cdots\longrightarrow
W_1\xrightarrow{\mu_1}W_0
\]
be a tower of tangent adapted toroidal blow-ups, and let $E\subset W_N$ be
the exceptional prime divisor created at the last step or the strict transform
of one created earlier.  Suppose that the tangential datum of $E$ is read on a
normalised invariant branch stratum $V_N\subset\sS_{W_N}^{\sn}$.  Let $F$ be
the corresponding divisor over the initial branch adjunction pair
$(V_0,B_{V_0})$.  Then
\[
\AF(E;W_0,\cG_0,\Delta_{W_0})=a(F;V_0,B_{V_0}).
\]
Equivalently, the equality between the foliated tangent discrepancy
contribution and the branch log-discrepancy contribution is additive along
tangent adapted towers.
\end{Lemma}

\begin{proof}
For a single tangent adapted blow-up, Proposition~\ref{prop:elementary-blowups}
computes the foliated log discrepancy in the foliated MMP convention.  Since
the exceptional divisor is invariant, the convention has $\epsilon(E)=0$;
therefore invariant components through the centre contribute with the same
coefficient-one rule as the boundary components appearing in foliated
adjunction, while transverse boundary components subtract their coefficients.
This gives exactly the ordinary log-discrepancy contribution on the
corresponding normalised branch pair.

Now pull the equality through the tower.  The transformation rule for
foliated canonical divisors under invariant blow-ups, in the convention used
in the foliated MMP \cite{CSMMP,SS,CSAdj}, gives for the composite
$\mu:W_N\to W_0$ the identity
\[
K_{\cG_N}+\Delta_{W_N}
=
\mu^*(K_{\cG_0}+\Delta_{W_0})
+
\sum_G a_G^{\mathcal F}G,
\]
where the sum runs over the exceptional invariant divisors created in the
tower and the coefficient of a fixed divisor is the sum of the relative
coefficients accumulated at the intermediate steps.  Equivalently,
$a_G^{\mathcal F}$ is the coefficient which defines the foliated log
discrepancy of $G$ over $(W_0,\cG_0,\Delta_{W_0})$ in this convention.
On the branch side, Theorem~\ref{thm:normalised-adjunction} gives the
analogous crepant pull-back formula
\[
K_{V_N}+B_{V_N}=\rho^*(K_{V_0}+B_{V_0})+A_\rho,
\]
where the coefficient of the divisor representing $E$ is the sum of the
ordinary branch contributions along the same sequence of centres.  The
single-step equality therefore adds term by term.  Hence the coefficient that
defines $\AF(E;W_0,\cG_0,\Delta_{W_0})$ equals the ordinary log discrepancy
$a(F;V_0,B_{V_0})$.
\end{proof}

\begin{Corollary}[Comparison for adapted toroidal invariant divisors]
\label{cor:foliated-comparison-tower}
Let $E$ be an invariant prime divisor obtained by tangent adapted toroidal
blow-ups.  Assume that its tangential datum is read on a normalised branch
stratum $V=S_\alpha^\nu$.  If $F$ is the divisor over $V$ representing $E$,
then
\[
\AF(E;X,\cF,\Delta)=a(F;V,B_V)=a_{\tan}(E;X,\cF,\Delta).
\]
\end{Corollary}

\begin{proof}
Choose an adapted toroidal model on which $E$ appears and resolve the branch
representative so that $F$ is visible.  The morphism from the initial adapted
model to this one is a tower of tangent adapted blow-ups.  By
Lemma~\ref{lem:foliated-discrepancy-additivity}, the foliated log discrepancy
of $E$ along this tower equals the ordinary log discrepancy of the
corresponding divisor on the normalised branch adjunction pair.  By definition,
$a_{\tan}$ is this ordinary branch log discrepancy for adapted toroidal branch
data.  The displayed equality follows.
\end{proof}

\begin{Remark}
The comparison is not asserted for arbitrary divisorial valuations over $X$.
It applies to adapted toroidal invariant divisors whose data are read on the
branch part of the separatrix system.  Conductor strata remain essential for
$\tmldtor$, but they are adjunction strata rather than prime divisors over
the foliated threefold.
\end{Remark}

\section{Tangential cylinders and the formula of Ein--Musta\c{t}\u{a}--Yasuda}
\label{sec:emy}

\begin{Definition}
Let $E$ be an adapted toroidal tangential divisorial datum represented by a divisor $F$ over a stratum $(V,B_V)$ of the separatrix system.  For $q\ge 1$, define $N_q^{\tan}(E)$ to be the tangential cylinder whose branch or conductor representative is the ordinary maximal divisorial cylinder $N_q(F)\subset J_\infty(V)$.
\end{Definition}

\begin{Definition}
The tangential logarithmic codimension of $N_q^{\tan}(E)$ is
\[
\lcodim_{\tan}(N_q^{\tan}(E)):=\lcodim_{(V,B_V)}(N_q(F)).
\]
This is independent of the representative by Proposition~\ref{prop:atan-invariance} and crepant compatibility.
\end{Definition}

\begin{Theorem}
\label{thm:tangential-emy}
For every adapted toroidal tangential divisorial datum $E$ and every $q\ge 1$,
\[
\lcodim_{\tan}\bigl(N_q^{\tan}(E)\bigr)=q\,a_{\tan}(E;X,\cF,\Delta).
\]
\end{Theorem}

\begin{proof}
Let \(E\) be an adapted toroidal tangential divisorial datum. By the branch--conductor classification, \(E\) is represented by an ordinary divisorial datum \(F\) over a normalised stratum \(V\), where \(V\) is either a normalised branch \(S_\alpha^\nu\) or a normalised conductor \(C_{\alpha\beta}^\nu\). Let \(B_V\) denote the corresponding branch or conductor boundary on \(V\). Thus the tangential maximal divisorial cylinder \(N_q^{\tan}(E)\) is represented, under the branch--conductor identification, by the ordinary maximal divisorial cylinder \(N_q(F)\subset J_\infty(V)\).
By definition of tangential logarithmic codimension through the branch--conductor system, one has
\[
\lcodim_{\tan} N_q^{\tan}(E)
=
\lcodim_{(V,B_V)} N_q(F).
\]
In the adapted toroidal charts considered here, the pair $(V,B_V)$ is log smooth at the generic point of the divisorial datum under consideration, or is replaced by a log resolution without changing the corresponding log discrepancy.  Thus the ordinary formula of Ein--Musta\c{t}\u{a}--Yasuda \cite{EMY} applies to the pair $(V,B_V)$ and gives
\[
\lcodim_{(V,B_V)} N_q(F)
=
q\,a(F;V,B_V).
\]
By the definition of tangential discrepancy and by the model independence proved in Proposition~\ref{prop:atan-invariance}, the ordinary discrepancy of the representing datum agrees with the adapted toroidal tangential discrepancy:
\[
a(F;V,B_V)
=
a_{\tan}(E;X,\cF,\Delta).
\]
Combining these equalities yields
\[
\lcodim_{\tan} N_q^{\tan}(E)
=
q\,a_{\tan}(E;X,\cF,\Delta).
\]
This proves the asserted formula.
\end{proof}

\begin{Definition}
For a closed tangent set $Y\subset X$, let $\mathcal E_{\mathrm tor}(Y)$ be the set of adapted toroidal tangential divisorial data $E$ such that $c_X(E)\subset Y$. Define
\[
\tmldtor(Y;X,\cF,\Delta)
:=
\inf_{E\in\mathcal E_{\mathrm tor}(Y)}
a_{\tan}(E;X,\cF,\Delta).
\]
Thus $\tmldtor$ is the toroidal/adapted invariant constructed in this article, not the infimum over all divisorial valuations over $X$.
\end{Definition}

\section{Detection of tangential log canonicity on the branch--conductor system}
\label{sec:ioa}

Let $\mathcal S_{\mathrm{tor}}(X,\cF)$ denote the set of normalised branch surfaces and normalised conductor curves occurring in the adapted toroidal separatrix system.  For $V\in\mathcal S_{\mathrm{tor}}(X,\cF)$ let $B_V$ denote its adjunction boundary.

\begin{Theorem}
\label{thm:ioa}
For every closed tangent set $Y$,
\[
\tmldtor(Y;X,\cF,\Delta)=
\inf_{V\in\mathcal S_{\mathrm{tor}}(X,\cF)}\mld(Y_V;V,B_V),
\]
where $Y_V$ is the inverse image of $Y$ on $V$, and conductor strata are counted once.
\end{Theorem}
 \begin{proof}
By definition,
\[
\tmldtor(Y;X,\cF,\Delta)
=
\inf_{E\in\mathcal E_{\mathrm{tor}}(Y)}
a_{\tan}(E;X,\cF,\Delta),
\]
where \(\mathcal E_{\mathrm{tor}}(Y)\) denotes the adapted toroidal
tangential divisorial data with centre contained in \(Y\).
Let \(E\in\mathcal E_{\mathrm{tor}}(Y)\).  By
Theorem~\ref{thm:classification}, after replacing the adapted model by a
common adapted toroidal refinement, \(E\) is represented by an ordinary
divisorial datum \(F\) over a normalised stratum
\[
V\in \Strata(W),
\qquad
V=S_\alpha^\nu
\quad\text{or}\quad
V=C_{\alpha\beta}^\nu .
\]
Let \(B_V\) be the adjunction boundary on \(V\).  By
Proposition~\ref{prop:atan-invariance},
\[
a_{\tan}(E;X,\cF,\Delta)
=
a(F;V,B_V).
\]
Thus
\[
\tmldtor(Y;X,\cF,\Delta)
=
\inf_{\substack{
(V,B_V)\\
F\ \mathrm{divisorial\ over}\ V\\
c_X(F)\subset Y
}}
a(F;V,B_V),
\]
where \(V\) ranges over the normalised branches and normalised conductors
appearing on adapted toroidal refinements of the separatrix system.

Conversely, let \(F\) be a divisorial datum over a normalised stratum
\(V\).  If \(V=S_\alpha^\nu\), the maximal divisorial set
\(N_q(F)\subset J_\infty(S_\alpha^\nu)\) defines an adapted tangential
divisorial cylinder by the branch-sector identification.  If
\(V=C_{\alpha\beta}^\nu\), the same construction gives a conductor-type
tangential divisorial cylinder through the coequaliser presentation of the
branch--conductor arc functor.  Hence every datum appearing on the right-hand
side comes from an adapted toroidal tangential divisorial datum on the left.

It remains to check multiplicities along conductors.  If
\(V=C_{\alpha\beta}^\nu\), the same conductor cylinder has two branch
representatives, one on \(S_\alpha^\nu\) and one on \(S_\beta^\nu\).  By
Proposition~\ref{prop:no-double-count}, these two representatives induce the
same adjunction boundary on \(C_{\alpha\beta}^\nu\), and the conductor datum
is counted once in the branch--conductor coequaliser.  Therefore no extra
factor or duplicate contribution appears in the infimum.
Combining the two inclusions gives
\[
\tmldtor(Y;X,\cF,\Delta)
=
\inf_{\substack{
V\in\mathcal S_{\mathrm{tor}}(X,\cF)\\
F\ \mathrm{divisorial\ over}\ V\\
c_X(F)\subset Y
}}
a(F;V,B_V),
\]
which is the claimed branch--conductor formula.
\end{proof}

\begin{Corollary}
The pair is tangentially log canonical along $Y$ if and only if each normalised branch and conductor pair $(V,B_V)$ is log canonical along $Y_V$.
\end{Corollary}


\section{Applications of the tangential formula}
\label{sec:applications}

Several formal consequences are recorded of the tangential codimension formula.
They should be viewed as arc-space refinements, in the adapted toroidal sector,
of adjunction phenomena known from the foliated MMP \cite{CSMMP,Spicer,SS,CSAdj}.  Throughout this section
\(\mathcal S_{\mathrm{tor}}(X,\cF)\) denotes the finite normalised
branch--conductor system associated with the fixed adapted toroidal category.
For \(V\in \mathcal S_{\mathrm{tor}}(X,\cF)\), write \((V,B_V)\) for the
ordinary adjunction pair and \(\rho_V:V\to X\) for the natural morphism.

\begin{Corollary}[Tangential toroidal inversion of adjunction]
\label{cor:tangential-ioa}
Let \(Y\subset X\) be a closed tangent set.  Under the hypotheses of the
tangential codimension formula,
\[
\tmldtor(Y;X,\cF,\Delta)
=
\inf_{V\in\mathcal S_{\mathrm{tor}}(X,\cF)}\mld(Y_V;V,B_V).
\]
In particular, \((X,\cF,\Delta)\) is tangentially toroidal log canonical
along \(Y\) if and only if every adjunction pair \((V,B_V)\) is log canonical
along \(Y_V\).
\end{Corollary}

\begin{proof}
This is Theorem~\ref{thm:ioa} together with the definition of tangential
toroidal log canonicity as non-negativity of the adapted toroidal tangential
minimal log discrepancy.  Equivalently, by Corollary~\ref{thm:intro-emy} and
its proof in Section~\ref{sec:emy}, every adapted toroidal tangential
divisorial cylinder is represented by an ordinary maximal divisorial cylinder
on a unique normalised branch or conductor stratum, up to the conductor
identifications, and its logarithmic codimension is the corresponding ordinary
log discrepancy.  Taking the infimum gives the displayed equality and the log
canonical criterion.
\end{proof}

\begin{Definition}
The adapted toroidal tangential non-lc and non-klt loci are
\[
\Nlc_{\tan,\mathrm{tor}}(X,\cF,\Delta)
:=\{p\mid \tmldtor(p;X,\cF,\Delta)<0\}
\]
and
\[
\Nklt_{\tan,\mathrm{tor}}(X,\cF,\Delta)
:=\{p\mid \tmldtor(p;X,\cF,\Delta)\le 0\}.
\]
The inequalities are understood for the adapted toroidal tangential
invariant constructed in this article.
\end{Definition}

\begin{Corollary}[Tangential non-lc and non-klt loci]
\label{cor:loci}
With notation as above,
\[
\Nlc_{\tan,\mathrm{tor}}(X,\cF,\Delta)
=
\bigcup_{V\in\mathcal S_{\mathrm{tor}}(X,\cF)}
\rho_V\bigl(\Nlc(V,B_V)\bigr).
\]
Similarly,
\[
\Nklt_{\tan,\mathrm{tor}}(X,\cF,\Delta)
=
\bigcup_{V\in\mathcal S_{\mathrm{tor}}(X,\cF)}
\rho_V\bigl(\Nklt(V,B_V)\bigr).
\]
Conductor strata are counted once.
\end{Corollary}

\begin{proof}
For a point \(p\) in the adapted tangential centre, the branch--conductor
formula gives
\[
\tmldtor(p;X,\cF,\Delta)
=
\min_{V\in\mathcal S_{\mathrm{tor}}(X,\cF)}
\inf_{q\in \rho_V^{-1}(p)}
\mld(q;V,B_V),
\]
with the convention that the infimum over an empty fibre is \(+\infty\).
By definition,
\[
p\in \Nlc_{\tan,\mathrm{tor}}(X,\cF,\Delta)
\]
if and only if
\[
\tmldtor(p;X,\cF,\Delta)<0.
\]
Using the displayed formula, this is equivalent to the existence of a stratum
\(V\in\mathcal S_{\mathrm{tor}}(X,\cF)\) and a point
\(q\in \rho_V^{-1}(p)\) such that
\[
\mld(q;V,B_V)<0.
\]
Equivalently,
\[
q\in \Nlc(V,B_V)
\quad\text{and}\quad
p=\rho_V(q).
\]
Thus
\[
\Nlc_{\tan,\mathrm{tor}}(X,\cF,\Delta)
=
\bigcup_{V\in\mathcal S_{\mathrm{tor}}(X,\cF)}
\rho_V\bigl(\Nlc(V,B_V)\bigr).
\]
The proof for the non-klt locus is identical, replacing the condition
\(\mld<0\) by the convention defining \(\Nklt\), equivalently
\(\mld\le 0\) in the present normalisation.  The convention that conductor
data are counted once is part of Theorem~\ref{thm:ioa}, so no additional
component is introduced by the two branch maps from a conductor.
\end{proof}

\begin{Lemma}[Lower semicontinuity on the branch--conductor strata]
\label{lem:lsc-strata}
Let \((V,B_V)\) be one of the normalised branch or conductor adjunction pairs
appearing in the adapted toroidal system.  Then the function
\[
q\longmapsto \mld(q;V,B_V)
\]
is lower semicontinuous on closed points of \(V\).
\end{Lemma}

\begin{proof}

Each stratum in the normalised branch--conductor system has dimension at most
two.  If \(V=C_{\alpha\beta}^{\nu}\) is a normalised conductor curve, then the
assertion follows from the one-dimensional formula for ordinary minimal log
discrepancies; at closed points it is given by \(1-\coeff_q(B_V)\).  If
\(V=S_\alpha^\nu\) is a normalised branch surface, it follows from Ambro's
two-dimensional lower semicontinuity theorem for ordinary minimal log
discrepancies \cite{AmbroMLD}.  Equivalently, after a log resolution of the
surface pair, the arc-space formula of Ein--Musta\c{t}\u{a}--Yasuda gives the
same statement on the smooth surface, and properness transfers it back to
\(V\) \cite{EMY}.
\end{proof}

\begin{Corollary}[Lower semicontinuity of the tangential toroidal mld]
\label{cor:lsc}
The function
\[
p\longmapsto \tmldtor(p;X,\cF,\Delta)
\]
is lower semicontinuous on the adapted tangential centre.
\end{Corollary}

\begin{proof}
Fix an object of the adapted toroidal category and let
\(\mathscr S_{\mathrm{tor}}\) be its finite branch--conductor system.  For each
\(V\in\mathscr S_{\mathrm{tor}}\), let \(\rho_V:V\to X\) be the natural proper
morphism and define
\[
\varphi_V(p):=
\inf_{q\in \rho_V^{-1}(p)}
\mld(q;V,B_V),
\]
with the convention that \(\varphi_V(p)=+\infty\) if the fibre is empty.  By
Lemma~\ref{lem:lsc-strata}, the mld function on \(V\) is lower
semicontinuous.  Therefore, for every real number \(a\), the set
\[
\{q\in V\mid \mld(q;V,B_V)\le a\}
\]
is closed, and its image under the proper map \(\rho_V\) is closed.  It
follows that \(\varphi_V\) is lower semicontinuous.  Since
\[
\tmldtor(p;X,\cF,\Delta)
=
\min_{V\in\mathscr S_{\mathrm{tor}}}\varphi_V(p),
\]
and the branch--conductor system is finite, \(\tmldtor\) is lower
semicontinuous.
\end{proof}

\begin{Corollary} 
\label{cor:cylinder-criterion}
The pair \((X,\cF,\Delta)\) is tangentially toroidal log canonical along a
closed tangent set \(Y\) if and only if
\[
a_{\tan}(E;X,\cF,\Delta)\ge 0
\]
for every adapted toroidal tangential divisorial datum \(E\) centred in \(Y\).
Equivalently, for every associated cylinder \(N_q^{\tan}(E)\),
\[
\lcodim_{\tan}\bigl(N_q^{\tan}(E)\bigr)\ge 0
\]
after the normalisation of logarithmic codimension used in Definition
\ref{def:tangential-cylinder-operational}.
\end{Corollary}

\begin{proof}
The first condition is the definition of tangential toroidal log canonicity in
terms of the adapted toroidal tangential discrepancies.  The equivalence with
logarithmic codimensions follows from the tangential
Ein--Musta\c{t}\u{a}--Yasuda formula
\[
\lcodim_{\tan}\bigl(N_q^{\tan}(E)\bigr)
= q\,a_{\tan}(E;X,\cF,\Delta),
\]
with \(q>0\).
\end{proof}

\begin{Remark}[Comparison with the foliated MMP approach]
Foliated adjunction and foliated inversion of adjunction are known by MMP
methods, especially through F-dlt modifications and the adjunction theory of
Cascini--Spicer and Spicer--Svaldi \cite{CSMMP,SS,CSAdj}.  Carter's jet
schemes of foliations give the infinitesimal language of strong tangency and
total separatrices \cite{Carter,CarterThesis}.  The corollaries above express
the corresponding tangential toroidal adjunction phenomena as equalities of
arc-space codimensions.  
\end{Remark}

\section{Adapted tangent/transverse decomposition}
\label{sec:transverse}

\begin{Proposition}
\label{prop:snc-comparison}
On an adapted model there are two ambient pairs which play different roles:
\[
B_W^{\tan}=\Dinv+\Delta_W,
\qquad
B_W^{\mathrm tr}=\Delta_W.
\]
For an adapted divisorial cylinder whose generic arc lies in $\Dinv$, the calculation of Ein--Musta\c{t}\u{a}--Yasuda \cite{EMY} for $(W,B_W^{\tan})$ restricts to the tangential computation on the corresponding branch or conductor.  For an adapted divisorial cylinder whose generic arc lies in $W\setminus \Dinv$, the invariant divisor is not part of the boundary and the computation is the ordinary one for the transverse ambient pair $(W,\Delta_W)$, hence after descent for $(X,\Delta)$.
\end{Proposition}

\begin{proof}
Let \(C\) be an irreducible cylinder and let \(\alpha_C:\Spec K[[t]]\to W\) be its generic arc. Suppose first that \(\alpha_C(\eta)\in \Dinv\), where \(\eta\) is the generic point of \(\Spec K[[t]]\). By Corollary~\ref{cor:branch-decomp}, \(\alpha_C\) is represented on a normalised stratum \(V\) of the branch--conductor system, where \(V\) is either a normalised branch \(S_\alpha^\nu\) or a normalised conductor \(C_{\alpha\beta}^\nu\). Let \(B_V\) be the corresponding boundary supplied by Theorem~\ref{thm:normalised-adjunction}.
By invariant-divisor adjunction, the restriction of the ambient logarithmic class to \(V\) is
\[
(K_{\mathcal G}+\Delta_W)|_V\sim_{\mathbb Q} K_V+B_V.
\]
Thus the logarithmic codimension used in the tangential sector agrees with the ordinary logarithmic codimension of the cylinder on the pair \((V,B_V)\). Applying the formula of Ein--Musta\c{t}\u{a}--Yasuda \cite{EMY} on \((V,B_V)\) gives precisely the tangential discrepancy computed by the branch--conductor system.

Now, suppose that \(\alpha_C(\eta)\notin \Dinv\). Then the arc may have finite contact order with \(\Dinv\), but it is not generically contained in the invariant separatrix divisor. Hence it is not part of the reduced tangential sector. In this transverse sector the invariant divisor \(\Dinv\) is not part of the boundary controlling tangential logarithmic codimension. The relevant ordinary pair is therefore the ambient pair with boundary \(\Delta_W\), not \(\Dinv+\Delta_W\).
Consequently the calculation of Ein--Musta\c{t}\u{a}--Yasuda \cite{EMY} gives the discrepancy of the ambient transverse pair \((W,\Delta_W)\). After descent from the adapted model to \(X\), this is the ordinary discrepancy contribution of the transverse sector. Thus the tangential branch--conductor computation applies precisely to the arcs generically contained in \(\Dinv\), while arcs generically outside \(\Dinv\) are computed by the ambient ordinary pair with boundary \(\Delta_W\).
\end{proof}

\begin{Corollary}
After passing to irreducible components, adapted toroidal divisorial cylinders decompose into a closed tangential sector governed by $\tmldtor$ and an open transverse sector governed by the ordinary invariant $\mld(X,\Delta)$.
\end{Corollary}



\section{Tangential Mather--Jacobian refinement}
\label{sec:mj}

In this section boundaries on possibly non-normal strata are always translated into the ideal-theoretic language used in Mather--Jacobian theory.  If
$
B_V=\sum_j b_jB_j
$
is a boundary on a reduced equidimensional stratum $V$, write
\begin{equation}\label{eq:mj-boundary-ideal}
\bbound_V:=\prod_j I_{B_j}^{b_j}
\end{equation}
as a formal product of ideals.  Thus
\[
a_{\MJ}(F;V,B_V)
\quad\text{means}\quad
a_{\MJ}(F;V,\bbound_V)
=\ord_F(\widehat K_{Y/V}-J_{Y/V})-\sum_j b_j\val_F(I_{B_j})+1.
\]

The Mather--Jacobian correction is attached to a singular separatrix space on
$X$, not to a smooth resolved branch where the correction is usually trivial.
It is therefore necessary to distinguish the resolved branch system on an adapted model from
the canonical seminormal image in $X$ associated with the fixed toroidal
category.


\begin{Construction}[Canonical image separatrix system on $X$]
\label{const:downstairs-sep}
Let $W\in\Adm(X,\cF,\Delta;W_0)$ and consider
\[
\sS_W^{\sn}\longrightarrow W\longrightarrow X.
\]
The reduced scheme-theoretic image of this morphism has a seminormalisation,
denoted $\sS_X^{\sn}$.  Its surface components are precisely the images of
those invariant separatrix branches whose image in $X$ has dimension two;
model-exceptional branches mapping to curves or points are retained only as
resolving strata and do not create new surface components on $X$.  The
normalisations of the surface components and of the reduced conductor images
are denoted $V_X^\nu$.  This construction uses only the fixed adapted
toroidal category and is algebraic, since it is formed from proper images of
algebraic divisors and their reduced intersections.
\end{Construction}

 


\begin{Proposition}
\label{prop:downstairs-canonical}
The germ of $\sS_X^{\sn}$ along a tangential centre is independent, up to unique isomorphism over $X$, of the object of $\Adm(X,\cF,\Delta;W_0)$ used in Construction \ref{const:downstairs-sep}, once the initial adapted model $W_0$ has been fixed.  More precisely, if two objects of this toroidal category are dominated by a common adapted toroidal refinement, then the reduced images of their seminormal separatrix spaces agree near the generic points of the corresponding tangential centres, and their seminormalisations are canonically identified.
\end{Proposition}

\begin{proof}
Let \(W_1,W_2\in\Adm(X,\cF,\Delta;W_0)\), and let
$
W_3 \xrightarrow{\mu_1} W_1
 $, $
W_3 \xrightarrow{\mu_2} W_2
$
be a common adapted toroidal refinement. Write
$
\pi_a:W_a\longrightarrow X,
 $ with $ a=1,2,3.
$
For each \(a\), let
\[
\mathfrak D_a=
\left(
\bigsqcup_{\alpha<\beta} C_{\alpha\beta,a}^{\nu}
\rightrightarrows
\bigsqcup_{\alpha} S_{\alpha,a}^{\nu}
\right)
\]
be the normalised branch--conductor diagram of \(W_a\). By Theorem~\ref{thm:common-refinement} and Proposition~\ref{prop:branch-functoriality}, there are morphisms of diagrams
\[
\mathfrak D_3\longrightarrow \mathfrak D_1,
\qquad
\mathfrak D_3\longrightarrow \mathfrak D_2.
\]
Equivalently, for \(a=1,2\), there are commutative diagrams
\[
\begin{tikzcd}
\displaystyle\bigsqcup_{\gamma<\delta} C_{\gamma\delta,3}^{\nu}
\arrow[r,shift left=.6ex] \arrow[r,shift right=.6ex] \arrow[d] &
\displaystyle\bigsqcup_{\gamma} S_{\gamma,3}^{\nu}
\arrow[d]\\
\displaystyle\bigsqcup_{\alpha<\beta} C_{\alpha\beta,a}^{\nu}
\arrow[r,shift left=.6ex] \arrow[r,shift right=.6ex] &
\displaystyle\bigsqcup_{\alpha} S_{\alpha,a}^{\nu}.
\end{tikzcd}
\]

After
shrinking \(X\) near the chosen tangential centre, the non-exceptional branches of the fixed initial object $W_0$ define a fixed finite
reduced collection of algebraic surface germs
\[
\mathcal S_X=\{S_1^X,\ldots,S_N^X\}
\]
such that every non-exceptional separatrix branch on every \(W_a\) maps to
one of the \(S_i^X\), and every \(S_i^X\) is obtained from a non-exceptional branch of the fixed initial model.  This follows because all objects of \(\Adm(X,\cF,\Delta;W_0)\) are obtained from \(W_0\) by coordinate stratum blow-ups, whose newly created exceptional branches have image of dimension at most one on \(X\). More
precisely, for every non-exceptional branch \(S_{\alpha,a}\subset W_a\)
there exists a unique index \(i(\alpha)\) such that
$
\overline{\pi_a(S_{\alpha,a})}_{\red}=S_{i(\alpha)}^X.
$
Exceptional branches created by the toroidal refinement are assigned no new
surface component on $X$; their images satisfy
\[
\overline{\pi_a(S_{\alpha,a})}_{\red}\subset S_i^X
\quad\text{or}\quad
\overline{\pi_a(S_{\alpha,a})}_{\red}\subset S_i^X\cap S_j^X
\]
for some \(i,j\).
Define the conductor curves on $X$ by scheme-theoretic intersection:
\[
C_{ij}^X:=(S_i^X\cap S_j^X)_{\red}.
\]
Then, for every conductor \(C_{\alpha\beta,a}\subset W_a\) whose adjacent
non-exceptional branches on $X$ are \(S_i^X\) and \(S_j^X\), one has
$
\overline{\pi_a(C_{\alpha\beta,a})}_{\red}=C_{ij}^X
$
near the generic point of the tangential centre. If one of the adjacent
branches is exceptional, then
$
\overline{\pi_a(C_{\alpha\beta,a})}_{\red}\subset C_{ij}^X
$
for the corresponding conductor image on $X$.

Hence the reduced image on $X$ of the whole branch--conductor diagram is
the fixed reduced scheme
\[
Z_X
:=
\left(
\bigcup_i S_i^X
\right)_{\red},
\qquad
\operatorname{Cond}(Z_X)
=
\left(
\bigcup_{i<j} C_{ij}^X
\right)_{\red}.
\]
In particular, for \(a=1,2,3\),
\[
\left(\pi_a\left(\bigcup_\alpha S_{\alpha,a}\right)\right)_{\red}=Z_X
\]
and
\[
\left(\pi_a\left(\bigcup_{\alpha<\beta}C_{\alpha\beta,a}\right)\right)_{\red}
=
\operatorname{Cond}(Z_X)
\]
after shrinking near the relevant generic point.

Let
$
Z_X^{\sn}\longrightarrow Z_X
$
be the seminormalisation. Since the reduced image \(Z_X\) on \(X\) is the
same for \(W_1\) and \(W_2\), the seminormalised bases over \(X\) obtained
from \(W_1\) and \(W_2\) are both canonically isomorphic to \(Z_X^{\sn}\).
Equivalently, there is a canonical identification over \(X\):
\[
\mathscr S_{X,1}^{\sn}
\cong
Z_X^{\sn}
\cong
\mathscr S_{X,2}^{\sn}.
\]

Thus the scheme on $X$ on which the Mather--Jacobian correction is
computed is independent of the adapted toroidal resolving model inside the
fixed category.  The Mather--Jacobian correction is therefore computed on the
fixed algebraic seminormal scheme \(Z_X^{\sn}\), while the upstairs
branch--conductor models serve only to resolve and present its strata.
\end{proof}

\begin{Definition}[Tangential Mather--Jacobian discrepancy]
Let $E$ be an adapted toroidal tangential divisorial datum represented by an ordinary divisorial datum $F$ over a fixed algebraic separatrix or conductor stratum $(V_X,\bbound_{V_X})$ on $X$.  Define
\[
a^{\tan}_{\MJ}(E;X,\cF,\Delta)=a_{\MJ}(F;V_X,\bbound_{V_X}),
\]
where the right-hand side is the ordinary Mather--Jacobian discrepancy of the scheme $V_X$ with the formal ideal product $\bbound_{V_X}$ encoding the boundary.
\end{Definition}

\begin{Proposition}[Mather--Jacobian invariance]
\label{prop:mj-invariance}
The number $a^{\tan}_{\MJ}$ is independent of the adapted resolved model dominating the fixed separatrix component.  If two adjacent branches induce the same normalised conductor stratum, the conductor value is computed once on that curve.
\end{Proposition}

\begin{proof}
By Proposition~\ref{prop:downstairs-canonical}, the algebraic stratum on $X$ representing the datum is independent of the adapted toroidal model. Denote this fixed stratum by \(V_X\), and let \(B_{V_X}\) be the corresponding boundary, encoded as a formal product of ideals on \(V_X\).

Let \(F\) be the divisorial datum whose Mather--Jacobian discrepancy is being computed. Choose a log resolution
\[
\rho:Y\longrightarrow V_X
\]
of \((V_X,B_{V_X})\) which factors through the Nash blow-up of \(V_X\). The ordinary Mather--Jacobian discrepancy is
\[
a_{\mathrm MJ}(F;V_X,B_{V_X})
=
\ord_F(\widehat K_{Y/V_X}-J_{Y/V_X})
-
\ord_F(B_{V_X})
+
1,
\]
where \(\widehat K_{Y/V_X}\) is the Mather discrepancy divisor and \(J_{Y/V_X}\) is the divisor associated with the pullback of the Jacobian ideal of \(V_X\). This number is independent of the chosen log resolution factoring through the Nash blow-up.

Now let \(W_a\) be an adapted toroidal model whose branch--conductor system represents the same datum, and let \(V_a\) be the corresponding normalised branch or conductor stratum mapping to \(V_X\). By Proposition~\ref{prop:downstairs-canonical}, the morphism \(V_a\to X\) factors through the fixed stratum \(V_X\):
\[
\begin{tikzcd}
V_a \arrow[dr] \arrow[r] & V_X \arrow[d]\\
& X.
\end{tikzcd}
\]
After replacing \(V_a\) by a further adapted toroidal refinement, there exists a common resolution
\[
\begin{tikzcd}
& Y' \arrow[dl] \arrow[dr] &\\
V_a \arrow[dr] && Y \arrow[dl,"\rho"]\\
& V_X &
\end{tikzcd}
\]
with \(Y'\to V_X\) factoring through the Nash blow-up of \(V_X\). Therefore the Mather and Jacobian discrepancy divisors pull back to \(Y'\), and the coefficient along the divisor representing \(F\) is the same as the coefficient computed on \(Y\). Hence
\[
a_{\mathrm MJ}^{\mathrm adapted}(F;V_a)
=
a_{\mathrm MJ}(F;V_X,B_{V_X}).
\]
Thus the Mather--Jacobian correction obtained from any adapted branch presentation is the ordinary Mather--Jacobian value \cite{dFD,IshiiMather,EI} on the fixed stratum \(V_X\) over \(X\).

In the conductor case, Proposition~\ref{prop:no-double-count} identifies the two adjacent branch restrictions with a single conductor boundary on the normalised conductor \(C_{\alpha\beta}^{\nu}\). Downstairs this gives a single fixed conductor stratum \(V_X=C_{ij}^X\). Since a normal curve over \(\mathbb C\) is smooth, the Mather discrepancy term on \(V_X\) is the ordinary curve discrepancy term, while the Jacobian correction is the usual correction attached to the chosen ideal product \(B_{V_X}\). Consequently the conductor contribution is computed once on \(V_X\), not once from each adjacent branch. This proves that the adapted Mather--Jacobian value is independent of the adapted toroidal model and of the chosen branch--conductor presentation.
\end{proof}

\begin{Definition}
Let an adapted toroidal tangential datum $E$ be represented by an ordinary divisorial datum $F$ over $V_X$. For the corresponding cylinder $N_q(F)\subset J_\infty(V_X)$, set
\[
\lcodim^{\MJ}_{\tan}(N_q^{\tan}(E))=
\lcodim^{\MJ}_{(V_X,\bbound_{V_X})}(N_q(F)).
\]
\end{Definition}

\begin{Theorem}
\label{thm:mj-cylinder}
For every adapted toroidal tangential divisorial datum $E$ represented by an ordinary divisorial datum $F$ over a fixed separatrix or conductor stratum, and every $q\ge1$,
\[
\lcodim^{\MJ}_{\tan}\bigl(N_q^{\tan}(E)\bigr)
=q\,a^{\tan}_{\MJ}(E;X,\cF,\Delta).
\]
\end{Theorem}

\begin{proof}
By Proposition~\ref{prop:downstairs-canonical}, the Mather--Jacobian calculation is performed on a fixed algebraic stratum \(V_X\), endowed with the formal product of ideals \(\bbound_{V_X}\). Let \(E\) be an adapted toroidal tangential divisorial datum and let \(F\) be the corresponding ordinary divisorial datum over \(V_X\). Under the branch--conductor identification, the tangential cylinder \(N_q^{\tan}(E)\) is identified with the ordinary maximal divisorial cylinder \(N_q(F)\subset J_\infty(V_X)\).

The ordinary Mather--Jacobian arc formula of de Fernex--Docampo, Ishii, and Ein--Ishii \cite{dFD,IshiiMather,EI} for the pair \((V_X,\bbound_{V_X})\) gives
\[
\codim_{J_\infty(V_X)}
N_q(F)
=
q\,a_{\mathrm MJ}(F;V_X,\bbound_{V_X}).
\]
By definition of the adapted tangential Mather--Jacobian discrepancy,
\[
a^{\tan}_{\MJ}(E;X,\cF,\Delta)
=
a_{\mathrm MJ}(F;V_X,\bbound_{V_X}).
\]
Moreover, by the branch--conductor identification of cylinders,
\[
\lcodim^{\MJ}_{\tan}\bigl(N_q^{\tan}(E)\bigr)
=
\codim_{J_\infty(V_X)}N_q(F).
\]
Combining these equalities gives
\[
\lcodim^{\MJ}_{\tan}\bigl(N_q^{\tan}(E)\bigr)
=
q\,a^{\tan}_{\MJ}(E;X,\cF,\Delta).
\]

Thus no new Mather--Jacobian arc formula is proved at the foliated level. The only additional ingredient is the branch--conductor presentation of the reduced tangential sector and the identification of the fixed stratum \(V_X\) over \(X\); after these identifications, the equality is precisely the ordinary Mather--Jacobian arc formula \cite{dFD,IshiiMather,EI} applied to \((V_X,\bbound_{V_X})\).
\end{proof}

\begin{Definition}
For a closed tangent set $Y\subset X$, let $\mathcal E_{\mathrm MJ}(Y)$ be the set of adapted toroidal tangential divisorial data $E$ with $c_X(E)\subset Y$ and with a fixed separatrix or conductor stratum on $X$.  Define
\[
\operatorname{tmld}_{\MJ}(Y;X,\cF,\Delta)
:=
\inf_{E\in\mathcal E_{\mathrm MJ}(Y)}
a^{\tan}_{\MJ}(E;X,\cF,\Delta).
\]
\end{Definition}

\begin{Theorem}
\[
\operatorname{tmld}_{\MJ}(Y;X,\cF,\Delta)=
\inf_{V_X\in\mathcal S_X(X,\cF)}\mld_{\MJ}(Y_{V_X};V_X,\bbound_{V_X}).
\]
\end{Theorem}

\begin{proof}
The proof is the valuation-theoretic infimum argument applied to the fixed Mather--Jacobian strata on $X$.
Let \(Y\subset X\) be a closed tangent set and let \(\mathcal E_{\mathrm MJ}(Y)\) be the set of adapted toroidal tangential divisorial data with centre contained in \(Y\) and with a fixed separatrix or conductor stratum on \(X\). By definition,
$
\operatorname{tmld}_{\MJ}(Y;X,\cF,\Delta)
=
\inf_{E\in\mathcal E_{\mathrm MJ}(Y)}
a^{\tan}_{\MJ}(E;X,\cF,\Delta).
$
For each \(E\in\mathcal E_{\mathrm MJ}(Y)\), Proposition~\ref{prop:downstairs-canonical} gives a fixed algebraic stratum \(V_X\), endowed with its formal product of ideals \(\bbound_{V_X}\), and Proposition~\ref{prop:mj-invariance} identifies the adapted tangential Mather--Jacobian discrepancy with the ordinary value on the fixed stratum:
\[
a^{\tan}_{\MJ}(E;X,\cF,\Delta)
=
a_{\mathrm MJ}(F;V_X,\bbound_{V_X}),
\]
where \(F\) is the ordinary divisorial datum over \(V_X\) corresponding to \(E\).
Therefore
\[
\operatorname{tmld}_{\MJ}(Y;X,\cF,\Delta)
=
\inf_{\substack{(V_X,\bbound_{V_X})\\ F\ \text{over }V_X\\ c_X(F)\subset Y}}
a_{\mathrm MJ}(F;V_X,\bbound_{V_X}),
\]
where the infimum is taken over the fixed Mather--Jacobian strata on $X$ arising from the canonical image branch--conductor system.
This is the same infimum argument as in Theorem~\ref{thm:ioa}, with ordinary log discrepancies replaced by ordinary Mather--Jacobian discrepancies on the strata on $X$. The model independence required to pass between adapted toroidal resolutions is exactly Proposition~\ref{prop:mj-invariance}.
\end{proof}

\subsection{Intrinsic comparison under algebraisation}
\label{subsec:intrinsic-mj-algebraisation}

The Mather--Jacobian refinement constructed above is relative to the
canonical image separatrix system associated with the adapted toroidal
category.  Under an additional algebraisation hypothesis, it becomes
intrinsic downstairs.

\begin{Assumption}[Algebraic separatrix hypothesis]
\label{ass:algebraic-separatrix}
Let $Y\subset X$ be the tangential centre under consideration.  Assume that
every intrinsic formal separatrix of $(X,\cF)$ relevant to $Y$ is
algebraisable, and that the resulting algebraic separatrix system, after
seminormalisation and normalisation of branches and conductors, coincides
with the canonical image separatrix system constructed from the adapted
toroidal category.
\end{Assumption}

\begin{Corollary}[Intrinsic Mather--Jacobian comparison under algebraisation]
\label{cor:intrinsic-mj-algebraisation}
Under Assumption~\ref{ass:algebraic-separatrix},
$a_{\MJ}^{\tan}(E;X,\cF,\Delta)$ agrees with the Mather--Jacobian
discrepancy of the intrinsic downstairs separatrix system.  In particular,
it is independent of the adapted image presentation.
\end{Corollary}

\begin{proof}
By definition, the tangential Mather--Jacobian discrepancy is computed on
the canonical image separatrix system.  If $E$ is represented by a divisor
$F$ over a stratum $V_X$, and the induced boundary is encoded by a formal
product of ideals $\bbound_{V_X}$, then
\[
a_{\MJ}^{\tan}(E;X,\cF,\Delta)
=
a_{\MJ}(F;V_X,\bbound_{V_X}).
\]
Under the algebraic separatrix hypothesis, the canonical image system agrees
with the intrinsic algebraic separatrix system downstairs.  Therefore
$V_X$, $F$, and $\bbound_{V_X}$ are the same objects in both constructions.
The displayed number is exactly the intrinsic downstairs Mather--Jacobian
discrepancy.  If a different adapted presentation is chosen, the same
hypothesis identifies it with the same intrinsic separatrix system and the
same boundary ideal.  Hence the value is independent of the adapted image
presentation.
\end{proof}

\subsection{Multiplier ideals on the seminormal separatrix space}

\begin{Construction}[Tangential Mather--Jacobian multiplier ideal]
Let $\rho_\alpha:S_\alpha^\nu\to\sS_X^{\sn}$ be the finite branch maps and let $\mathfrak a$ be an ideal on $X$ whose pull-back is non-zero on every branch.  Define $\cJ^{\tan}_{\MJ}(\mathfrak a^t)$ as the subsheaf of $\cO_{\sS_X^{\sn}}$ whose pull-back to the normalisation $\coprod S_\alpha^\nu$ is
\[
\left\{(s_\alpha)_\alpha\in\bigoplus_\alpha
\cJ_{\MJ}(S_\alpha^\nu,\bbound_{S_\alpha};(\mathfrak a\cO_{S_\alpha^\nu})^t)
:\ s_\alpha|_{C_{\alpha\beta}^\nu}=s_\beta|_{C_{\alpha\beta}^\nu}\right\}.
\]
Equivalently it is the equaliser inside $\cO_{\sS_X^{\sn}}$ of the two conductor restriction maps.
The restrictions in the equaliser are the images of the inclusions
\[
\cJ_{\MJ}(S_\alpha^\nu,\bbound_{S_\alpha};(\mathfrak a\cO_{S_\alpha^\nu})^t)\subset \cO_{S_\alpha^\nu}
\]
under the quotient maps \(\cO_{S_\alpha^\nu}\to\cO_{C_{\alpha\beta}^\nu}\).  No equality of the restricted ideals from the two sides is assumed a priori; the equaliser imposes equality of their images as sections of \(\cO_{C_{\alpha\beta}^\nu}\).
\end{Construction}


The equaliser may be written explicitly as
\begin{equation}\label{eq:mj-equaliser-ideal}
\cJ^{\tan}_{\MJ}(\mathfrak a^t)=
\left\{(s_\alpha)_\alpha\in\bigoplus_\alpha \cJ_{\MJ,\alpha}:
\mathrm{res}_{\alpha\beta}^{\alpha}(s_\alpha)=
\mathrm{res}_{\alpha\beta}^{\beta}(s_\beta)\text{ for all }\alpha<\beta
\right\},
\end{equation}
inside the equaliser presentation of $\cO_{\sS_X^{\sn}}$.  The maps to $\cO_{C_{\alpha\beta}^\nu}$ are obtained by composing the inclusions $\cJ_{\MJ,\alpha}\subset\cO_{S_\alpha^\nu}$ with the quotient maps $\cO_{S_\alpha^\nu}\to\cO_{C_{\alpha\beta}^\nu}$; no equality of the restricted ideals is assumed a priori.  In this construction
\[
\cJ_{\MJ,\alpha}=
\cJ_{\MJ}(S_\alpha^\nu,\bbound_{S_\alpha};(\mathfrak a\cO_{S_\alpha^\nu})^t).
\]

\begin{center}
\begin{tikzcd}[column sep=large]
\cJ^{\tan}_{\MJ}(\mathfrak a^t) \arrow[r,hook] \arrow[d,hook] &
\displaystyle\bigoplus_\alpha \cJ_{\MJ,\alpha}
\arrow[r,shift left=.55ex] \arrow[r,shift right=.55ex] \arrow[d,hook] &
\displaystyle\bigoplus_{\alpha<\beta}\cO_{C_{\alpha\beta}^\nu}
\arrow[d,equal] \\
\cO_{\sS_X^{\sn}} \arrow[r,hook] &
\displaystyle\bigoplus_\alpha \cO_{S_\alpha^\nu}
\arrow[r,shift left=.55ex] \arrow[r,shift right=.55ex] &
\displaystyle\bigoplus_{\alpha<\beta}\cO_{C_{\alpha\beta}^\nu}
\end{tikzcd}
\end{center}

\begin{Theorem}
\label{thm:coherence}
With the conductor-compatible equaliser definition above, the sheaf $\cJ^{\tan}_{\MJ}(\mathfrak a^t)$ is a coherent ideal sheaf on $\sS_X^{\sn}$.
\end{Theorem}

\begin{proof}
The normalisation map $\coprod S_\alpha^\nu\to\sS_X^{\sn}$ is finite.  Ordinary Mather--Jacobian multiplier ideals on the normal branches are coherent, and finite pushforward preserves coherence.  The conductor compatibility condition is an equaliser of coherent sheaves on the Noetherian seminormal scheme $\sS_X^{\sn}$.  Since the equaliser is taken inside the equaliser presentation of $\cO_{\sS_X^{\sn}}$, it is an ideal subsheaf of $\cO_{\sS_X^{\sn}}$.
\end{proof}

\begin{Definition}[Tangential Mather--Jacobian threshold]
For a closed tangent set $Y$, define
\[
\tlct_{\MJ}(Y;X,\cF,\Delta;\mathfrak a)=
\inf_{V_X\in\mathcal S_X(X,\cF)}\lct_{\MJ}(Y_{V_X};V_X,\bbound_{V_X};\mathfrak a\cO_{V_X}).
\]
\end{Definition}

\begin{Proposition}
\label{prop:threshold-normalisation}
The pull-back of $\cJ^{\tan}_{\MJ}(\mathfrak a^t)$ to the normalisation $ \bigsqcup S_\alpha^\nu$ is trivial over the inverse image of $Y$ for $t<\tlct_{\MJ}$ and is non-trivial there for $t>\tlct_{\MJ}$.  If, in addition, triviality of coherent ideals on $\sS_X^{\sn}$ is detected after the finite normalisation map in the given conductor neighbourhood, then the same statement holds on $\sS_X^{\sn}$ itself.
\end{Proposition}

\begin{proof}
Let
\[
\nu:\widetilde V_X:=\bigsqcup_i V_i\longrightarrow V_X
\]
be the finite normalisation map of the seminormal branch--conductor stratum on $X$, and write
\[
\nu^{-1}\bbound_{V_X}=\{\bbound_{V_i}\}_i
\]
for the induced formal products of ideals on the normalised branches. By construction of the branch--conductor system, coherent ideals on \(V_X\) are identified with tuples of coherent ideals on the \(V_i\) whose restrictions agree on the normalised conductor. Thus, for every real number \(c\geq 0\), the Mather--Jacobian multiplier ideal on \(V_X\) is recovered as the equaliser
\[
\mathcal J_{\mathrm MJ}(V_X,\bbound_{V_X}^{\,c})
=
\ker\left(
\bigoplus_i \mathcal J_{\mathrm MJ}(V_i,\bbound_{V_i}^{\,c})
\rightrightarrows
\bigoplus_{i<j}
\mathcal J_{\mathrm MJ}(C_{ij}^{\nu},\bbound_{C_{ij}}^{\,c})
\right),
\]
where the two arrows are the restriction maps from the adjacent normalised branches to the normalised conductor.
After pulling back by \(\nu\), the assertion becomes the ordinary Mather--Jacobian multiplier-ideal criterion \cite{EI,dFD} on each normal branch:
\[
\mathcal J_{\mathrm MJ}(V_i,\bbound_{V_i}^{\,c})=\mathcal O_{V_i}
\quad\Longleftrightarrow\quad
c<\operatorname{lct}_{\mathrm MJ}(V_i,\bbound_{V_i}).
\]
Therefore the first value of \(c\) for which the branch multiplier ideal becomes non-trivial is
\[
\inf_i \operatorname{lct}_{\mathrm MJ}(V_i,\bbound_{V_i}),
\]
which is precisely the infimum over the branch--conductor data appearing in the definition of \(\tlct_{\mathrm MJ}\).

Under the conductor-compatibility hypothesis imposed in the construction, the conductor equaliser does not introduce a smaller threshold: it only imposes equality of the already-defined branch sections after restriction to the normalised conductor. Hence the global ideal remains trivial precisely until one of the branch Mather--Jacobian multiplier ideals becomes non-trivial. Consequently,
\[
\tlct_{\mathrm MJ}(V_X,\bbound_{V_X})
=
\inf_i \operatorname{lct}_{\mathrm MJ}(V_i,\bbound_{V_i}).
\]
Finally, the descent statement is the elementary finite-normalisation criterion for coherent ideals. Namely, if \(\mathfrak a,\mathfrak b\subset\mathcal O_{V_X}\) are coherent ideals such that their pullbacks to \(\widetilde V_X\) agree and their conductor restrictions agree, then \(\mathfrak a=\mathfrak b\). Applying this to the two multiplier ideals in question proves the asserted equality on \(V_X\).
\end{proof}

\subsection{Conditional ordinary consequences}

The preceding Mather--Jacobian comparison transfers ordinary
singularity-theoretic consequences to the branch--conductor pairs.  We record
the consequences obtained by combining this comparison with the ordinary
Mather--Jacobian rationality and Du Bois criteria on the normalised strata,
together with Du Bois gluing for the seminormal pushout.

\begin{Assumption}[MJ rationality and Du Bois gluing hypotheses]
\label{ass:mj-db-gluing}
In a neighbourhood of the tangential set \(Y\), assume that the normalised
branch and conductor pairs lie in the range of the ordinary
Mather--Jacobian criteria of de Fernex--Docampo and Ein--Ishii
\cite{dFD,EI}.  Thus the relevant strata are reduced and equidimensional,
the Jacobian and boundary ideals define admissible pairs, and the
corresponding rationality or Du Bois criterion applies on each normalised
stratum.

Assume moreover that the seminormal branch--conductor pushout satisfies the
finite seminormal Du Bois gluing hypotheses of Koll\'ar--Kov\'acs
\cite{KollarKovacs}.  Equivalently, in the conductor neighbourhood under
consideration, the pushout is obtained from Du Bois branches by gluing along
Du Bois conductor strata.
\end{Assumption}

\begin{Theorem}
\label{thm:rational-dubois}
Under Assumption~\ref{ass:mj-db-gluing}, if every branch pair is
Mather--Jacobian canonical near \(Y\), then the normalised branches have
rational singularities near \(Y\).  Under the same gluing hypotheses, if
every branch pair is Mather--Jacobian log canonical near \(Y\), then
\(\sS_X^{\sn}\) is Du Bois near the inverse image of \(Y\).
\end{Theorem}

\begin{proof}
Let
\[
\nu:\widetilde{\sS}^{\sn}
=
\bigsqcup_i V_i
\longrightarrow
\sS^{\sn}
\]
be the finite normalisation of the seminormal separatrix space.  Its
conductor diagram is
\[
\begin{tikzcd}[column sep=large]
\displaystyle\bigsqcup_{i<j} C_{ij}^{\nu}
\arrow[r,shift left=.55ex]
\arrow[r,shift right=.55ex]
&
\displaystyle\bigsqcup_i V_i
\arrow[r]
&
\sS^{\sn}.
\end{tikzcd}
\]
For each normalised branch \(V_i\), let \(B_i\) denote the induced boundary,
or formal product of ideals, coming from the tangential
Mather--Jacobian datum.  For a normalised conductor \(C_{ij}^{\nu}\), the
conductor-compatible construction of the data requires the two restrictions
of the idealistic boundary to define the same conductor datum.  We denote it
by
\[
B_{ij}:=B_i|_{C_{ij}^{\nu}}=B_j|_{C_{ij}^{\nu}}.
\]
Here Proposition~\ref{prop:no-double-count} supplies the corresponding
compatibility for the adjunction divisor classes, while equality of the
formal ideal products is part of the conductor-compatible Mather--Jacobian
input used in this section.
We first prove the rational assertion.  This assertion is branchwise.  For
each \(i\), the ordinary Mather--Jacobian canonical hypothesis gives
\[
a_{\MJ}(F;V_i,B_i)\geq 1
\]
for every exceptional divisor \(F\) over \(V_i\) in the range of the ordinary
criterion.  By the Mather--Jacobian criterion of de Fernex--Docampo and
Ein--Ishii \cite{dFD,EI}, it follows that $V_i$ has rational singularities near the inverse image of the centre.  Hence
\[
\widetilde{\sS}^{\sn}
=
\bigsqcup_i V_i
\]
has rational singularities branchwise.
We now prove the Du Bois assertion.  For each \(i\), the ordinary
Mather--Jacobian log canonical hypothesis gives
\[
a_{\MJ}(F;V_i,B_i)\geq 0
\]
for every divisor \(F\) over \(V_i\) in the range of the ordinary criterion.
Again by the ordinary Mather--Jacobian criterion, $V_i$ is Du Bois.
Moreover, each \(C_{ij}^{\nu}\) is a normalised curve; over \(\bC\) it is
regular, hence Du Bois.  Thus
\[
V_i \text{ is Du Bois},\qquad
C_{ij}^{\nu} \text{ is Du Bois},\qquad
B_i|_{C_{ij}^{\nu}}=B_j|_{C_{ij}^{\nu}}.
\]
The space \(\sS^{\sn}\) is the seminormal pushout of the conductor diagram
above.  By the gluing hypothesis in the statement, the finite seminormal
Du Bois gluing theorem of Koll\'ar--Kov\'acs applies to the union
\[
\sS^{\sn}
=
\left(\bigsqcup_i V_i\right)
/\!\sim_{\bigsqcup C_{ij}^{\nu}}
\]
along the Du Bois conductor strata.  Consequently, \(\sS^{\sn}\) is
Du Bois near the image of the centre.  Thus rational singularities are
obtained on the normalised branches, while
the Du Bois property descends to the seminormal separatrix space by
conductor-compatible Du Bois gluing.
\end{proof}

\begin{Proposition}
If every branch or conductor stratum on $X$ is a local complete intersection and the Jacobian correction is trivial along the datum considered, then, for every adapted toroidal tangential datum $E$ represented by a divisor $F$ over the corresponding stratum,
\[
a^{\tan}_{\MJ}(E;X,\cF,\Delta)=a_{\tan}(E;X,\cF,\Delta).
\]
Consequently $\operatorname{tmld}_{\MJ}=\tmldtor$ for the adapted toroidal data to which both sides apply.
\end{Proposition}

\begin{proof}
Assume that the stratum $V$ on $X$ is a local complete intersection along the centre of the datum and that the Jacobian correction is trivial there. Let \(B_V\) be the corresponding branch or conductor boundary, encoded as a formal product of ideals. Since \(V\) is a local complete intersection, the Mather discrepancy divisor and the Jacobian discrepancy divisor cancel in the Mather--Jacobian discrepancy formula:
\[
\widehat K_{Y/V}-J_{Y/V}=K_{Y/V}
\]
on any log resolution \(Y\to V\) dominating the Nash blow-up. Hence, for every divisor \(F\) over \(V\),
\[
a_{\mathrm MJ}(F;V,B_V)=a(F;V,B_V).
\]
Now, let \(E\) be an adapted toroidal tangential divisorial datum represented by \(F\) over the same normalised branch or conductor pair \((V,B_V)\). By definition of the two tangential discrepancies,
\[
a^{\tan}_{\MJ}(E;X,\cF,\Delta)
=
a_{\mathrm MJ}(F;V,B_V)
\]
and
$
a_{\tan}(E;X,\cF,\Delta)
=
a(F;V,B_V).
$
Therefore
\[
a^{\tan}_{\MJ}(E;X,\cF,\Delta)
=
a_{\tan}(E;X,\cF,\Delta).
\]
Thus, in the locally complete intersection range with trivial Jacobian correction, the Mather--Jacobian tangential discrepancy and the ordinary toroidal tangential discrepancy evaluate the same representative divisor over the same branch or conductor pair and hence agree.
\end{proof}

\section{Local refinements and birational comparison}
\label{sec:expanded-verification}

This section records the local verifications used previously in a form intended to specify where the foliation enters the argument.  The statements are elementary in SNC coordinates, but they are the points at which the branch--conductor theory differs from a direct application of the ordinary theory to a divisor.

\subsection{Double and triple separatrix points}

\begin{Lemma}
Let $p$ be a point at which two invariant branches meet, say $S_1=\{x=0\}$ and $S_2=\{y=0\}$.  In the completed local ring the tangential formal ideal is $(xy)$.  Hence a non-constant irreducible tangential cylinder centred at $p$ is either generically on $S_1$, generically on $S_2$, or generically on the conductor curve $\{x=y=0\}$.
\end{Lemma}

\begin{proof}
By Theorem~\ref{thm:formal-rep}, the reduced tangential arc functor centred on \(\Sigma_{\tan}\) is identified with the reduced arc functor of the invariant divisor \(\Dinv\). In the present local chart, \(\Dinv=(xy=0)\). Hence a \(K[[t]]\)-arc \(\gamma:\Spec K[[t]]\to W\) in the reduced tangential sector satisfies
\[
x(\gamma)y(\gamma)=0
\]
in \(K[[t]]\). Since \(K[[t]]\) is an integral domain, this implies
$
x(\gamma)=0
 $ or $
y(\gamma)=0.$
Thus every reduced tangential arc factors through at least one of the two invariant branches \(S_x=(x=0)\) or \(S_y=(y=0)\).
If the generic point of \(\gamma\) lies on exactly one of these branches, then \(\gamma\) belongs to the corresponding branch sector. If the generic point satisfies both equations, then it factors through the conductor curve
$
C_{xy}=S_x\cap S_y=(x=y=0).
$
Therefore the reduced tangential arc sector is the union of the two branch sectors, glued along the conductor sector.
For an irreducible cylinder \(C\), apply the same argument to its generic arc \(\gamma_C:\Spec K(C)[[t]]\to W\). The generic arc determines whether \(C\) is generically supported on \(S_x\), on \(S_y\), or on the conductor \(C_{xy}\). Hence the branch--conductor decomposition holds componentwise for irreducible cylinders.
\end{proof}

\begin{Lemma}
Let $p$ be a point at which three invariant branches $x=0$, $y=0$, $z=0$ meet.  Every adapted toroidal tangential divisorial datum centred at $p$ becomes branch or conductor type after a finite sequence of adapted blow-ups of coordinate strata.
\end{Lemma}

\begin{proof}
By definition, an adapted toroidal datum at the triple point is represented by a primitive vector
\[
v=(a,b,c)\in\bN^3,
\qquad \gcd(a,b,c)=1,
\]
in the cone $\sigma=\bR_{\ge0}^3$ of the coordinate arrangement $xyz=0$.  The associated monomial valuation is
\[
\nu_v(x^iy^jz^k)=ai+bj+ck.
\]
A star subdivision of $\sigma$ along the ray $\bR_{\ge0}v$ produces a toroidal blow-up whose exceptional divisor has valuation $\nu_v$.  The blow-up of the origin is the first such subdivision when $v=(1,1,1)$; in the $x$-chart it has coordinates
$
y=xu, z=xv,
$
and the reduced transform is $xuv=0$.  For a general primitive vector, perform the Euclidean sequence of star subdivisions which inserts the ray $\bR_{\ge0}v$ into the fan.  This sequence is finite because each star subdivision refines a rational polyhedral cone and the desired ray is fixed.

On the resulting toroidal model the centre of $\nu_v$ is the divisor corresponding to the inserted ray.  Its generic point lies either on a single invariant component of the refined divisor or on a pairwise intersection with another component, according to the cone containing the adjacent rays.  These two cases are precisely the branch and conductor presentations.  No assertion is made for non-toroidal divisorial valuations centred at the triple point.
\end{proof}

\subsection{One-step transform formula}

\begin{Proposition}
\label{prop:one-step-full}
Let $S$ be an invariant branch with boundary $\Theta_S$, and let $p\in S$ be a tangent point centre lying on boundary components of total coefficient $b$.  Let $\sigma:S'\to S$ be the blow-up of the smooth surface point $p$, and let $F$ be the exceptional curve.  Then the crepant transform of the branch boundary is
\[
\Theta_{S'}^{\mathrm cr}=\sigma^{-1}_*\Theta_S+(b-1)F,
\]
and
$
K_{S'}+\Theta_{S'}^{\mathrm cr}=\sigma^*(K_S+\Theta_S).
$
The ordinary log discrepancy of $F$ with respect to $(S,\Theta_S)$ is $2-b$.
\end{Proposition}

\begin{proof}
For the blow-up of a smooth surface point one has
$
K_{S'}=\sigma^*K_S+F.
$
If $D_1,\ldots,D_\ell$ are the components of $\Theta_S$ through $p$, with coefficients $b_1,\ldots,b_\ell$, and $b=\sum_i b_i$, then
\[
\sigma^*\Theta_S=\sigma^{-1}_*\Theta_S+bF.
\]
Consequently
$
\sigma^*(K_S+\Theta_S)
=K_{S'}+\sigma^{-1}_*\Theta_S+(b-1)F.
$
This is precisely the stated crepant identity.
The log discrepancy of the exceptional curve is computed from
\[
K_{S'}+\sigma^{-1}_*\Theta_S
=\sigma^*(K_S+\Theta_S)+(1-b)F,
\]
so that $a(F;S,\Theta_S)=1+(1-b)=2-b$.  If the ambient blow-up is along a conductor curve, its restriction to an old branch is the blow-up of a Cartier divisor and is therefore an isomorphism; the ambient exceptional divisor is then a new branch in the refined branch--conductor system, not an exceptional curve on the old branch surface.
\end{proof}

\begin{Corollary}
\label{cor:iterated-branch-discrepancy}
Let
\[
W_N\longrightarrow W_{N-1}\longrightarrow\cdots\longrightarrow W_0=W
\]
be a tower of adapted tangent point blow-ups over a normalised branch
\(S_\alpha^\nu\).  For every exceptional curve \(F_k\) which appears on the
corresponding branch surface, the ordinary log discrepancy of \(F_k\) with
respect to the transformed branch pair is equal to the tangential discrepancy
of the associated adapted tangent exceptional divisor \(E_k\):
\[
a(F_k;S_{\alpha,k}^{\nu},B_{\alpha,k})
=
a_{\tan}(E_k;X,\cF,\Delta).
\]
\end{Corollary}

\begin{proof}
Let \(F_k\) be an exceptional curve in the branch tower, and let \(k\) be the
first index at which it appears.  By Proposition~\ref{prop:one-step-full},
the one-step adapted tangent blow-up gives
\[
a(F_k;S_{\alpha,k}^{\nu},B_{\alpha,k})
=
a_{\tan}(E_k;X,\cF,\Delta).
\]
For \(m\geq k\), subsequent adapted tangent blow-ups only replace the pair by
a higher birational model.  The discrepancy of the fixed divisorial valuation
defined by \(F_k\) is invariant under further birational pull-back:
\[
a(F_k;S_{\alpha,k}^{\nu},B_{\alpha,k})
=
a(F_k;S_{\alpha,m}^{\nu},B_{\alpha,m}).
\]
The associated adapted exceptional divisor \(E_k\) defines the corresponding
tangential divisorial datum on the threefold model, and the foliated
adjunction comparison identifies its tangential discrepancy with the same
ordinary branch discrepancy.  Hence
\[
a(F_k;S_{\alpha,m}^{\nu},B_{\alpha,m})
=
a_{\tan}(E_k;X,\cF,\Delta)
\]
for every \(m\geq k\).  This proves the claim for every exceptional curve in
the tower.
\end{proof}


\section{Jet-tangent interpretation of the Mather--Jacobian correction}
\label{sec:jet-tangent-mj}

The Mather--Jacobian correction can also be read from tangent spaces to jet schemes of the branch strata on $X$.  This is not an additional construction; it is the ordinary jet-theoretic formula for Jacobian discrepancies of de Fernex--Docampo \cite{dFD} applied after the reduction to the branch--conductor system.

\begin{Proposition}
\label{prop:jet-tangent}
Let $E$ be an adapted toroidal tangential divisorial datum represented by an ordinary divisor $F$ over an $n$-dimensional stratum $V_X$ on $X$.  For $m\gg0$, the Jacobian discrepancy part of $a^{\tan}_{\MJ}(E;X,\cF,\Delta)$ is computed by the tangent space to the $m$-jet scheme of $V_X$ at the generic jet associated with $F$:
\[
k_F^{\diamond}(V_X)+1
=
2n(m+1)-\dim T(V_X)_m|_{\eta_{F,m}},
\]
with the ordinary normalisation of de Fernex--Docampo \cite{dFD}.  Adding the boundary contribution of $B_{V_X}$ gives the formula for the pair $(V_X,B_{V_X})$.
\end{Proposition}

\begin{proof}
The displayed equality is the ordinary jet--tangent formula for the Jacobian discrepancy on the fixed scheme on $X$ \(V_X\). More precisely, let \(F\) be the ordinary divisorial datum over \(V_X\) corresponding to the adapted toroidal tangential datum \(E\). Choose a resolution
$
\rho:Y\longrightarrow V_X
$
which factors through the Nash blow-up and on which \(F\) appears as a divisor. The ordinary Mather--Jacobian formula gives
\[
a_{\mathrm MJ}(F;V_X,\bbound_{V_X})
=
\ord_F(\widehat K_{Y/V_X}-J_{Y/V_X})
-
\ord_F(\bbound_{V_X})
+
1.
\]
Equivalently, in jet-theoretic terms, the codimension of the ordinary divisorial cylinder \(N_q(F)\subset J_\infty(V_X)\) is
\[
\codim_{J_\infty(V_X)}N_q(F)
=
q\,a_{\mathrm MJ}(F;V_X,\bbound_{V_X}).
\]

By the branch--conductor construction on $X$, the tangential cylinder \(N_q^{\tan}(E)\) is identified with the ordinary cylinder \(N_q(F)\) on \(V_X\). Hence
\[
\lcodim^{\MJ}_{\tan}\bigl(N_q^{\tan}(E)\bigr)
=
\codim_{J_\infty(V_X)}N_q(F).
\]
Using the definition
$
a^{\tan}_{\MJ}(E;X,\cF,\Delta)
=
a_{\mathrm MJ}(F;V_X,\bbound_{V_X}),
$
one obtains
\[
\lcodim^{\MJ}_{\tan}\bigl(N_q^{\tan}(E)\bigr)
=
q\,a^{\tan}_{\MJ}(E;X,\cF,\Delta).
\]
Thus the tangential statement contains no additional jet-theoretic assertion beyond the ordinary Mather--Jacobian formula on \(V_X\); the sole additional ingredient is the identification of the reduced tangential cylinder with the ordinary cylinder on the fixed stratum on $X$.
\end{proof}

\section{Examples and model computations}
\label{sec:examples}

\begin{Example}
Let $W=\mathbb A^3$ with coordinates $x,y,z$, and let
$
\Dinv=S_x+S_y,
S_x=(x=0),
S_y=(y=0).$
The conductor is
$
C_{xy}=S_x\cap S_y=(x=y=0).
$
Assume that the foliation is logarithmic simple adapted along $\Dinv$,
for instance locally generated, up to a unit and saturation, by a form whose
logarithmic part is
\[
\lambda_x\frac{dx}{x}+\lambda_y\frac{dy}{y},
\qquad
\lambda_x,\lambda_y\ne 0,
\qquad
\lambda_x+\lambda_y\ne 0.
\]
By invariant-divisor adjunction, the boundary on the branch $S_x$ is
$
\Theta_x=C_{xy}.$
Let $u=y|_{S_x}$ be a local equation of $C_{xy}$ on $S_x\simeq\mathbb A^2$.
For $q\geq 1$, consider the ordinary contact cylinder
\[
C_q=\{\alpha\in J_\infty(S_x):\ord_t(u\circ\alpha)=q\}.
\]
Since $u$ is a smooth coordinate on the surface $S_x$, the ordinary
codimension of $C_q$ in $J_\infty(S_x)$ is $q$. The boundary contribution is
also $q$, because $C_{xy}$ appears in $\Theta_x$ with coefficient one.
Therefore
\[
\lcodim_{(S_x,\Theta_x)}(C_q)=q-q=0.
\]
Thus contact with the other invariant separatrix contributes no positive
tangential logarithmic codimension. This is the local model for the coefficient-one
cancellation corresponding to the foliated tangent convention
$\epsilon(E)=0$.
Equivalently, blow up the conductor curve $(x=y=0)$. On the chart $y=xu$,
the transformed invariant divisor has components
$
E=(x=0)$ and $
S_y'=(u=0).
$ The logarithmic residues transform as
\[
\lambda_x\frac{dx}{x}+\lambda_y\frac{dy}{y}
=
(\lambda_x+\lambda_y)\frac{dx}{x}
+
\lambda_y\frac{du}{u}.
\]
The exceptional divisor $E$ is invariant because
$\lambda_x+\lambda_y\ne 0$. On the transformed branch, the trace of $E$
appears with coefficient one, and the same cancellation occurs on the
refined branch pair.
\end{Example}

\begin{Example}
Let
$
\Dinv=S_x=(x=0)
$
and let
$
T=(z=0)
$
be a transverse boundary component with coefficient $a$, where $0\leq a\leq
1$. On the invariant branch $S_x$, the adjunction boundary is
$
\Theta_x=a\,T|_{S_x}.
$
Let $v=z|_{S_x}$. For $q\geq 1$, consider the cylinder
\[
C_q^T=\{\alpha\in J_\infty(S_x):\ord_t(v\circ\alpha)=q\}.
\]
Since $v=0$ is a smooth coordinate curve on the smooth surface $S_x$, the
ordinary codimension of $C_q^T$ is $q$. The boundary contribution is $aq$.
Hence
\[
\lcodim_{(S_x,\Theta_x)}(C_q^T)=q-aq=q(1-a).
\]
This number is the branch contribution of the tangent
extraction over the curve $S_x\cap T$.
If $a=0$, the transverse divisor is absent from the boundary and the
logarithmic codimension is $q$. If $a=1$, the transverse trace behaves, from the
point of view of logarithmic codimension, like a coefficient-one boundary component
and the contribution vanishes. Thus transverse components are measured by
the defect $1-a$, whereas invariant traces always occur with coefficient
one by foliated adjunction.
\end{Example}

\begin{Example}
Let
$
\Dinv=S_x+S_y,
S_x=(x=0),
S_y=(y=0),
$
and let $T=(z=0)$ be transverse with coefficient $a$. On $S_x$, the branch
boundary is
$
\Theta_x=C_{xy}+a\,T|_{S_x}.
$
Write
 $
u=y|_{S_x},$ and $
v=z|_{S_x}.
$
For integers $m,n\geq 0$, consider the cylinder
\[
C_{m,n}
=
\{\alpha\in J_\infty(S_x):\ord_t(u\circ\alpha)=m,\
\ord_t(v\circ\alpha)=n\}.
\]
The curves $u=0$ and $v=0$ are coordinate curves on the smooth surface
$S_x$. Therefore the imposed conditions are independent and
$
\codim_{J_\infty(S_x)}(C_{m,n})=m+n.
$
The boundary contribution is
$
\ord_{C_{m,n}}(\Theta_x)=m+an.
$
Thus
\[
\lcodim_{(S_x,\Theta_x)}(C_{m,n})
=
(m+n)-(m+an)
=
(1-a)n.
\]
The invariant contact order $m$ cancels, while the transverse contact order
$n$ remains weighted by $1-a$. This is the elementary local computation
behind the general branch--conductor logarithmic codimension formula.
\end{Example}


\begin{Example}
Continue with
$
\Dinv=(xyz=0).$
Consider the conductor
$
C_{xy}=S_x\cap S_y=(x=y=0).
$
On the branch $S_x$, the boundary contains the invariant traces $
C_{xy}
 $ and $
C_{xz}
$
with coefficient one. Thus, near the triple point,
$
\Theta_x=C_{xy}+C_{xz}.
$
Similarly, on the branch $S_y$,$
\Theta_y=C_{xy}+C_{yz}.
$
To compare the two restrictions along $C_{xy}$, remove the common conductor
component $C_{xy}$ and restrict the remaining boundary to
$C_{xy}^{\nu}$. Then
\[
C_{xz}|_{C_{xy}^{\nu}}=(x=y=z=0)
=
C_{yz}|_{C_{xy}^{\nu}}
\]
as reduced zero-dimensional divisors. Hence
$
(\Theta_x-C_{xy})|_{C_{xy}^{\nu}}
=
(\Theta_y-C_{xy})|_{C_{xy}^{\nu}}.
$
The triple-point contribution is therefore a single conductor gluing datum,
not two independent branch contributions. This is the local model for the
no-double-counting statement.
\end{Example}



\begin{Example}
Let
$
\Dinv=S_x+S_y,
T=(z=0),
\Delta_W=aT.
$
On $S_x$, the boundary is
\[
\Theta_x=C_{xy}+aT|_{S_x}.
\]
For a cylinder with
$
\ord(C_{xy})=m,
\ord(T|_{S_x})=n,
$
the ordinary codimension in $J_\infty(S_x)$ is
$
m+n.
$ However, the tangential logarithmic codimension is
$
(m+n)-(m+an)=(1-a)n.
$
Thus two cylinders with different values of $m$ but the same value of $n$
have different ordinary codimension but the same tangential logarithmic codimension.
This illustrates why ordinary codimension is not the birational invariant in
the tangential branch sector. The invariant quantity is the log-codimension
after subtracting the branch adjunction boundary.
\end{Example}

\begin{Example}
Let $V_X$ be a
singular separatrix stratum on $X$, and let $\bbound_{V_X}$ be the
formal product of ideals encoding the induced boundary. Let
$
\rho:Y\to V_X
$
be a log resolution factoring through the Nash blow-up. For a divisor $F$
on $Y$, the ordinary Mather--Jacobian discrepancy is
\[
a_{\mathrm MJ}(F;V_X,\bbound_{V_X})
=
\ord_F(\widehat K_{Y/V_X}-J_{Y/V_X})
-
\ord_F(\bbound_{V_X})
+
1.
\]
The term $J_{Y/V_X}$ records the Jacobian correction. Equivalently, it
detects the order of the Jacobian ideal of $V_X$ along $F$.
Suppose that $F_1$ and $F_2$ are two divisors over $V_X$ such that their
ordinary log discrepancies with respect to the same boundary agree:
$
a(F_1;V_X,\bbound_{V_X})=
a(F_2;V_X,\bbound_{V_X}),
$
but their Jacobian orders differ:
$
\ord_{F_1}(\jaco_{V_X})\ne \ord_{F_2}(\jaco_{V_X}).
$
Then the Mather--Jacobian discrepancies may differ:
$
a_{\mathrm MJ}(F_1;V_X,\bbound_{V_X})
\ne
a_{\mathrm MJ}(F_2;V_X,\bbound_{V_X}).
$
The corresponding adapted tangential Mather--Jacobian discrepancies are
obtained by representing the same data through the branch--conductor system:
\[
a^{\tan}_{\MJ}(E_i;X,\cF,\Delta)
=
a_{\mathrm MJ}(F_i;V_X,\bbound_{V_X}).
\]
Thus the Mather--Jacobian tangential refinement detects singularity defects
of the separatrix stratum on $X$ which are invisible to the ordinary
adapted toroidal tangential discrepancy.
\end{Example}

\begin{Example}[The locally complete intersection range]
Let \(V_X\) be an algebraic separatrix stratum on \(X\), and assume that
\(V_X\) is locally complete intersection near the centre of the datum.  Let
$
\mu:Y\longrightarrow V_X
$
be a log resolution of the pair \((V_X,\bbound_{V_X})\) which factors through
the Nash blow-up, and let \(F\) be a prime divisor on \(Y\).  For locally
complete intersection varieties one has the standard comparison
\[
\widehat K_{Y/V_X}-J_{Y/V_X}=K_{Y/V_X}.
\]
Taking coefficients along \(F\), this gives
$
\ord_F(\widehat K_{Y/V_X})
-
\ord_F(J_{Y/V_X})
=
\ord_F(K_{Y/V_X}).
$
Consequently the Mather--Jacobian discrepancy and the ordinary discrepancy
of the stratum pair coincide:
\[
a_{\MJ}(F;V_X,\bbound_{V_X})
=
a(F;V_X,\bbound_{V_X}).
\]

If an adapted toroidal tangential datum \(E\) is represented by \(F\) over
the stratum \(V_X\), then the tangential Mather--Jacobian discrepancy also
coincides with the ordinary tangential discrepancy:
\[
a^{\tan}_{\MJ}(E;X,\cF,\Delta)
=
a_{\tan}(E;X,\cF,\Delta).
\]
Thus, on locally complete intersection separatrix strata, the
Mather--Jacobian refinement introduces no additional Jacobian correction and
reduces to the ordinary adapted toroidal tangential discrepancy.
\end{Example}


\paragraph{\bf Acknowledgements.} The author is partially supported by the Universit\`a degli Studi di Bari and is a member of INdAM-GNSAGA.

\ \begin{thebibliography}{99}

\bibitem{AmbroMLD}
\bibaut{F. Ambro},
\emph{On minimal log discrepancies},
Math. Res. Lett. \textbf{6} (1999), no. 5, 573--580.

\bibitem{Can04}
\bibaut{F. Cano},
\emph{Reduction of the singularities of codimension one singular foliations
in dimension three},
Ann. of Math. (2) \textbf{160} (2004), no. 3, 907--1011.

\bibitem{Carter}
\bibaut{P. J. Carter},
\emph{On the resolutions of non-dicritical foliations},
arXiv:2203.11094, 2024.

\bibitem{CarterThesis}
\bibaut{P. J. Carter},
\emph{On the resolution of foliation singularities via jet schemes},
Ph.D. thesis, University of Liverpool, 2020.

\bibitem{CC92}
\bibaut{F. Cano and D. Cerveau},
\emph{Desingularization of nondicritical holomorphic foliations and
existence of separatrices},
Acta Math. \textbf{169} (1992), no. 1--2, 1--103.

\bibitem{CSMMP}
\bibaut{P. Cascini and C. Spicer},
\emph{MMP for co-rank one foliations on threefolds},
Invent. Math. \textbf{225} (2021), no. 2, 603--690.

\bibitem{CSRankOne}
\bibaut{P. Cascini and C. Spicer},
\emph{On the MMP for rank one foliations on threefolds},
Forum Math. Pi \textbf{13} (2025), e4,
arXiv:2012.11433.

\bibitem{CSAdj}
\bibaut{P. Cascini and C. Spicer},
\emph{Foliation adjunction},
Math. Ann. \textbf{391} (2025), 5695--5727.

\bibitem{dFD}
\bibaut{T. de Fernex and R. Docampo},
\emph{Jacobian discrepancies and rational singularities},
J. Eur. Math. Soc. \textbf{16} (2014), no. 1, 165--199.

\bibitem{EI}
\bibaut{L. Ein and S. Ishii},
\emph{Singularities with respect to Mather--Jacobian discrepancies},
arXiv:1310.6882, 2013.

\bibitem{EMY}
\bibaut{L. Ein, M. Musta\c{t}\u{a} and T. Yasuda},
\emph{Jet schemes, log discrepancies and inversion of adjunction},
Invent. Math. \textbf{153} (2003), no. 3, 519--535.

\bibitem{IshiiMather}
\bibaut{S. Ishii},
\emph{Mather discrepancy and the arc spaces},
Ann. Inst. Fourier \textbf{63} (2013), no. 1, 89--111.

\bibitem{KollarKovacs}
\bibaut{J. Koll\'ar and S. Kov\'acs},
\emph{Log canonical singularities are Du Bois},
J. Amer. Math. Soc. \textbf{23} (2010), no. 3, 791--813.

\bibitem{SS}
\bibaut{C. Spicer and R. Svaldi},
\emph{Local and global applications of the Minimal Model Program for co-rank one foliations on threefolds},
J. Eur. Math. Soc. \textbf{24} (2022), no. 11, 3969--4025.

\bibitem{Spicer}
\bibaut{C. Spicer},
\emph{Higher-dimensional foliated Mori theory},
Compos. Math. \textbf{156} (2020), no. 1, 1--38.

\end{thebibliography}

\end{document}
